\documentclass[11pt]{article}

% --- arxiv-style template (kourgeorge/arxiv-style) ---
\usepackage{arxiv}

% --- engine: xelatex (Unicode-friendly) ---
\usepackage{fontspec}
\setmainfont{Times New Roman}
\setmonofont{Menlo}[Scale=0.85]

% --- core packages ---
\usepackage{amsmath, amssymb}
\usepackage{booktabs}
\usepackage{longtable}
\usepackage{array}
\usepackage{xcolor}
\usepackage{hyperref}
\hypersetup{
    colorlinks=true,
    linkcolor=blue,
    citecolor=blue,
    urlcolor=blue,
}
\usepackage{graphicx}
\usepackage{microtype}
\usepackage{caption}

% --- url-aware line breaking ---
\usepackage{url}
\PassOptionsToPackage{hyphens}{url}
\Urlmuskip=0mu plus 1mu      % allow stretch around URL break-points
\sloppy                       % prefer loose lines over overflow
\setlength{\emergencystretch}{3em}

% --- inline-code (pandoc \passthrough) breakability ---
% Long file paths and identifiers that pandoc renders as
% \texttt{...} are otherwise unbreakable, which produced 200+pt
% overfull \hbox warnings. \seqsplit inserts a discretionary
% break opportunity between every character; LaTeX only breaks
% there when no better option exists, so short inline-code spans
% stay intact in the normal flow.
\usepackage{seqsplit}

% --- code listings ---
\usepackage{listings}
\lstset{
    basicstyle=\ttfamily\small,
    breaklines=true,
    columns=fullflexible,
    keepspaces=true,
}

% --- bibliography (bibtex, more portable than biber) ---
\usepackage[numbers]{natbib}

% --- pandoc compatibility helpers ---
\usepackage{calc}
\providecommand{\tightlist}{%
  \setlength{\itemsep}{0pt}\setlength{\parskip}{0pt}}
\providecommand{\pandocbounded}[1]{#1}
\providecommand{\real}[1]{#1}
% pandoc emits \texttt{...} (and sometimes \passthrough{...}) around
% inline-code spans. Override BOTH so long file paths and identifiers
% break across lines instead of producing 200+pt overfull \hbox
% warnings. \seqsplit only inserts break opportunities; LaTeX only
% breaks there when no better option exists, so short inline-code
% spans stay intact in the normal flow.
\newcommand{\passthrough}[1]{\texttt{\seqsplit{#1}}}
\let\OldTexttt\texttt
\renewcommand{\texttt}[1]{\OldTexttt{\seqsplit{#1}}}

% Long tables (Tables 11/12 with 7-column engine grids carry rows
% like "707 302 µs" + "not reported by runner" + "39.06 MB †" that
% otherwise overflow column widths). Render every longtable in
% \footnotesize so the wide rows fit; the surrounding prose stays
% in 11pt body size.
\usepackage{etoolbox}
\AtBeginEnvironment{longtable}{%
  \scriptsize%                   % was \footnotesize — clears the
                                  % last 1pt overflow in Table 12
  \setlength{\tabcolsep}{3pt}% tighter cell padding (default 6pt)
  \renewcommand{\arraystretch}{0.95}% tighter row spacing
}
% pandoc emits \def\LTcaptype{none} around longtables to suppress
% caption numbering; longtable then tries \addtocounter{\LTcaptype}.
% Define the dummy counter so tectonic does not bail.
\newcounter{none}

% --- title block ---
\title{Dazzle: An Embedded Database for On-Device LLM Agents}

\author{
    Ivan Aliaga \\
    Universidad Ricardo Palma \\
    Lima, Peru \\
    \texttt{ivan.aliaga@urp.edu.pe}
}

\date{April 2026}

% --- arxiv metadata ---
\renewcommand{\headeright}{Preprint}
\renewcommand{\undertitle}{Preprint}
\renewcommand{\shorttitle}{Dazzle: An Embedded Database for On-Device LLM Agents}

% --- begin document ---
\begin{document}

\maketitle

\begin{abstract}
Mobile applications that embed a language-model agent now need a
persistent state layer between inference calls --- sensor windows,
materialised aggregates, approximate counts, vector indices over
recent context. The embedded databases available on the platform
(SQLite, RocksDB, LMDB, ObjectBox) ship two to four data-structure
primitives each, so every higher-order construct (bounded streams,
per-field TTL, HyperLogLog, HNSW vector search) is reimplemented
in application code, on top of those primitives, by every team
that needs them.

We present \textbf{Dazzle}, a fork of Valkey (the Linux
Foundation's Redis-compatible store) whose server runs
\emph{inside} the process of Android and iOS applications rather
than as a daemon over TCP loopback. The fork ships ten native
data-structure primitives, including HNSW vector search, against
$\leq 4$ in any embedded competitor we measure on the same
platforms; it exposes them through a typed mobile SDK on Maven
Central, pub.dev, npm and Swift Package Manager.

We characterise the latency envelope across twelve backends on
two physical devices (Moto G35 5G, iPhone 12 Pro). Steady-state
retrieval on the materialised-aggregates path lands at 50~µs on
Moto G35 5G and 7.24~µs on iPhone 12 Pro; HNSW vector search
stays sub-millisecond at recall $\geq 0.95$ across the
9-configuration grid we measure. At 50~µs against a
$\sim 2.84$~s on-device LLM inference turn the retrieval cost is
\textbf{0.00176\,\%} of the inference call --- below the noise
floor of the end-to-end loop. The deciding axis for an on-device
agent backend is therefore \textbf{primitive surface and
developer ergonomics}, not microbenchmark deltas. We report the
absolute per-engine microsecond numbers anyway, because the same
engines also serve non-LLM workloads (analytics, sync, indexing)
where microseconds matter.

An end-to-end 2$\times$2 ablation on 200 Natural Questions queries
on the Moto G35 5G shows that the retrieval primitive is what
drives factual accuracy at this scale. Qwen 2.5 0.5B with Dazzle
RAG reaches 0.630 \texttt{EM\_contains} against 0.105 without
retrieval ($6.0\times$); Qwen 2.5 1.5B with RAG reaches 0.735
against 0.110 without ($6.7\times$). A 380~MB model with
retrieval beats a 940~MB model ($3\times$ larger) without
retrieval on every factual metric, end-to-end on a
\$150-class Android phone. The bottleneck is not the model size,
the SoC, or the microsecond cost of retrieval; it is whether the
backend lets the agent loop reach a vector index at all.
\end{abstract}

\keywords{embedded database \and on-device LLM \and mobile agents
  \and Valkey \and HNSW vector search \and Apache 2.0}

% --- body content (generated from paper_v2_en.md by pandoc) ---
\section{Introduction}\label{introduction}

Quantized language models under 3 GB now run at interactive rates on
smartphone-class ARM64 hardware: Gemma 2 \citep{gemma2} and Gemma 4
\citep{gemma4}; Phi-3 mini \citep{phi3}; the Llama 3 herd \citep{llama3}
including Llama 3.2 1B/3B; MobileLLM \citep{mobilellm}; and the Qwen 2.5 series
\citep{qwen25} used in §5.9. The runtime stack --- \texttt{llama.cpp}
\citep{llamacpp} with its GGUF quantised-weight format \citep{gguf}, LiteRT-LM
\citep{litertlm}, MLC-LLM \citep{mlcllm}, ExecuTorch \citep{executorch} --- is
mature and reaches sub-second decoded-token latency on \$150-class handsets.
What has not matured at the same pace is the \textbf{state layer} an agent needs
--- a place to hold what it has seen between inference calls, coordinate across
turns, and survive process restarts.

Consider a mobile IoT monitoring agent: it collects temperature, humidity,
pressure, and battery readings at intervals; an LLM processes batches at
checkpoints and produces synthesis reports. Between checkpoints the agent needs
(1) a time-ordered bounded buffer over the most recent readings, (2) running
aggregates (min/max/avg/ count) without re-scanning history, (3) probabilistic
counting for anomaly and cardinality signals, and (4) per-field expiration to
keep the context block within the prompt budget. Each requirement maps cleanly
to a specific data structure --- bounded streams, sorted sets and hashes,
HyperLogLog, TTL on hash fields. None of these are exotic. None are research
problems.

What is missing is a backend that offers them. On Android and iOS, the dominant
embedded options --- SQLite \citep{sqlite}, RocksDB \citep{rocksdb} (an LSM-tree
descendant of LevelDB \citep{leveldb}), LMDB \citep{lmdb}, ObjectBox
\citep{objectbox4} --- and the analytical embedded peer DuckDB \citep{duckdb}
--- expose between two and four data structures total: rows, key-value pairs,
blob stores, occasionally a vector index. Every higher-order construct (bounded
streams, materialized aggregates, approximate counts, TTL semantics, similarity
search) has to be reimplemented in application code, on top of those primitives,
by every team that needs them. The resulting code is fragile under concurrent
access, the atomicity guarantees that come for free with a mature in-process
server are silently lost, and the same handful of bugs is rediscovered in app
after app.

The latency math frees the argument from the usual microbenchmark race. Even the
slowest embedded backend in our evaluation --- SQLite at roughly 3 ms retrieval
at N = 20 000 under concurrent load on a \$150 Android handset --- sits three
orders of magnitude below a single on-device LLM inference turn (1--3 seconds
for a 1--3 B-parameter model). Optimizing microseconds of retrieval is not what
moves the product needle. What does is which primitives the backend exposes, how
ergonomic they are to invoke, and how much glue code the agent author has to
write before the state layer behaves.

We argue that the right axis of comparison for on-device agent backends is
\textbf{primitive surface and developer ergonomics}, not raw microbenchmark
throughput. We support this thesis with four pieces of evidence: (a) a
characterization of Dazzle's latency envelope showing it sits comfortably within
the LLM-noise floor across all configurations tested; (b) a
primitive-by-primitive comparison showing \textbf{ten native primitives in
Dazzle against \(\leq\) 4 in any embedded competitor on the same platforms}; (c)
a development-complexity comparison in lines of code for the same agent state
layer; and (d) an end-to-end demonstration that the primitives Dazzle exposes
enable qualitatively new patterns --- small-model + on-device
retrieval-augmented generation \citep{lewisrag, karpukhindpr} --- that are not
achievable when the only available primitive is a flat key-value or row store.

\textbf{Contributions.}

\begin{enumerate}
\def\labelenumi{\arabic{enumi}.}
\tightlist
\item
  \textbf{First in-process port of Valkey to Android and iOS application
  processes}: a self-pipe transport, write-through snapshot cache, and SPSC ring
  buffer on Android replace Valkey's TCP loopback while preserving the full
  primitive surface and Lua scripting.
\item
  \textbf{One measured engineering contribution plus two future-proofing layers}
  (§3, §5.4). The write-through snapshot cache is the single layer that moves
  throughput on the measured Write-Heavy / Read-Light workload --- the §5.4
  factorial 2³ ablation shows the parallel worker pool and the hash-bucket
  dispatch path are no-ops on this workload (within run-to-run noise) because
  the snapshot cache already absorbs the contention they were designed to
  remove. We retain those two layers in the architecture as keyspace-scaling
  headroom: under benchmarks not run in this paper (much larger keyspaces,
  multiple writer threads, fan-in from several producer queues) we expect them
  to earn back their cost, and the §5.4 numbers describe the conditions under
  which they do not. The post-EVAL auto-mirror that keeps Lua-mutated state
  visible to the cache, and the in-process transport itself, are the two further
  engineering contributions of independent interest documented in §3.
\item
  \textbf{Latency-envelope characterization} across twelve backends, with N from
  200 to 20 000, on a budget Android handset (Moto G35 5G, Unisoc T760) and a
  physical iPhone 12 Pro --- full bench-coverage parity (every backend runs on
  both platforms; iOS retains the Phase-0 self-pipe transport --- see §6.3 ---
  and the SPSC ring buffer / \texttt{io\_uring} optimisations are Android-only),
  including the ObjectBox 5.3 native iOS port wired in this revision (entities,
  Sourcery-generated bindings, and a simulator- runnable XCTest suite under
  \texttt{experiment/storage/ios/Tests/ObjectBoxTests.swift}).
\item
  \textbf{Primitive-surface and development-complexity comparison} against five
  embedded alternatives (SQLite, LMDB, RocksDB, ObjectBox, InMemory) plus a
  Valkey-over-TCP reference: ten native primitives versus \(\leq\) 4, with
  concrete LOC measurements for the same agent state layer (§5.6).
\item
  \textbf{End-to-end demonstration} via a 2×2 RAG ablation on 200 NQ queries
  (§5.9, Table 15). Adding the Dazzle vector primitive lifts
  \texttt{EM\_contains} 6.0× on Qwen 2.5 0.5B (0.105 → 0.630) and 6.7× on Qwen
  2.5 1.5B (0.110 → 0.735); the small-model-with- retrieval cell beats the
  large-model-without-retrieval cell on every factual metric. The result speaks
  to what becomes implementable when the right primitive is present, not to a
  microsecond delta over a competitor.
\item
  \textbf{Public SDK distribution} of \texttt{dazzle-sdk} 1.0.0-beta.4 across
  four standard registries (Maven Central, pub.dev, npm, Swift Package Manager),
  with five LLM adapters behind a common \texttt{LLMClient} interface (Appendix
  A). Apache-2.0 (Dazzle) / BSD-3-Clause (derivative Valkey portions).
\end{enumerate}

\section{Background and Motivation}\label{background-and-motivation}

\subsection{State Requirements for an On-Device LLM
Agent}\label{state-requirements-for-an-on-device-llm-agent}

An LLM-as-agent differs from single-turn inference in that it must:

\begin{itemize}
\tightlist
\item
  Maintain a bounded observation window (circular buffer with automatic
  eviction).
\item
  Track running aggregates (min/max/avg without O(N) re-scan).
\item
  Approximately count distinct events (HyperLogLog in 12 KB at
  \textasciitilde0.8 \% relative error \citep{hyperloglog} vs O(N) exact).
\item
  Apply per-field TTL to stale observations (automatic expiration without app
  cron).
\item
  Signal between components (Pub/Sub without polling).
\end{itemize}

These primitives have existed for a decade in Redis/Valkey. What is new is
embedding them \textbf{inside a mobile app process}, eliminating the network hop
between agent and store.

\subsection{Native Primitives by Backend}\label{native-primitives-by-backend}

\textbf{Table 1 --- Per-primitive availability per embedded backend.} Each cell
records \emph{how} the primitive is reached, not just \emph{whether} it is
reachable. \textbf{N} = native typed API call shipped by the engine itself.
\textbf{E} = official extension shipped by the upstream project (e.g.~SQLite's
R\emph{Tree, in the amalgamation since 2008). \textbf{A} = app-level
reimplementation --- the developer writes glue code over lower primitives in the
engine. \textbf{T} = third-party extension distributed separately from the
engine. The ``Total Native'' row counts only \textbf{N}; \textbf{E}, \textbf{A}
and \textbf{T} mean the primitive is }achievable* but with different engineering
cost, captured quantitatively by the LOC differential in Table 9 and discussed
explicitly in §6.3.

{\def\LTcaptype{none} % do not increment counter
\begin{longtable}[]{@{}
  >{\raggedright\arraybackslash}p{(\linewidth - 10\tabcolsep) * \real{0.4400}}
  >{\centering\arraybackslash}p{(\linewidth - 10\tabcolsep) * \real{0.1067}}
  >{\centering\arraybackslash}p{(\linewidth - 10\tabcolsep) * \real{0.1067}}
  >{\centering\arraybackslash}p{(\linewidth - 10\tabcolsep) * \real{0.1200}}
  >{\centering\arraybackslash}p{(\linewidth - 10\tabcolsep) * \real{0.1467}}
  >{\centering\arraybackslash}p{(\linewidth - 10\tabcolsep) * \real{0.0800}}@{}}
\toprule\noalign{}
\begin{minipage}[b]{\linewidth}\raggedright
Primitive
\end{minipage} & \begin{minipage}[b]{\linewidth}\centering
Dazzle
\end{minipage} & \begin{minipage}[b]{\linewidth}\centering
SQLite
\end{minipage} & \begin{minipage}[b]{\linewidth}\centering
RocksDB
\end{minipage} & \begin{minipage}[b]{\linewidth}\centering
ObjectBox
\end{minipage} & \begin{minipage}[b]{\linewidth}\centering
LMDB
\end{minipage} \\
\midrule\noalign{}
\endhead
\bottomrule\noalign{}
\endlastfoot
Bounded stream with MAXLEN & N & A & A & A & A \\
Sorted sets with range query & N & A¹ & A & N & A \\
Atomic float increment & N & A² & A & A & A \\
Per-field TTL & N & A & A & A & A \\
HyperLogLog & N & A & A & A & A \\
Geo indexing & N & E³ & A & A & A \\
Pub/Sub & N & A & A & N⁴ & A \\
Server-side scripting (Lua) & N & A & A & A & A \\
Vector search (HNSW) & N & T⁵ & A & N & A \\
Cursor-based iteration & N & N & N & N & N \\
\textbf{Total Native (N)} & \textbf{10} & \textbf{1} & \textbf{1} & \textbf{4} &
\textbf{1} \\
\end{longtable}
}

¹ Implementable via \texttt{ORDER\ BY\ …\ WHERE\ …\ BETWEEN} in app SQL. ²
Implementable via \texttt{UPDATE\ …\ SET\ v\ =\ v\ +\ ?} inside a transaction. ³
R*Tree is an official SQLite extension shipped inside the amalgamation since
2008. ⁴ ObjectBox \texttt{DataSubscription} is a typed 1-call API for
change-notification consumers; we count it as a native pub/sub equivalent for
the agent-loop use case (single consumer subscribed to a typed stream of
changes), even though it is observer-pattern rather than topic-routed. ⁵ SQLite
has no native or official vector primitive; vector search is provided only by
third-party extensions (commercial \texttt{sqliteai/sqlite-vector}, open-source
\texttt{asg017/sqlite-vec}), which we benchmark separately in §5.8 as commercial
peers rather than as a SQLite native feature.

The headline count (10 vs \(\leq\) 4 native) measures \textbf{how many
primitives the agent author can reach with a single typed library call} --- the
``primitive surface'' axis of the Dazzle thesis. Primitives marked \textbf{E},
\textbf{A} or \textbf{T} in non-Dazzle backends are \emph{implementable}; the
question is engineering cost, which we measure directly in §5.6 (Table 9: 175
LOC Android / 186 LOC iOS for Dazzle vs 200--290 LOC for the SQLite-based
equivalents). A reviewer reading Table 1 as ``primitives that exist anywhere in
this engine's ecosystem'' would arrive at a much smaller gap; that is a
different (and weaker) question than the one Table 1 is structured to answer,
and we make this axis-of-comparison choice explicit in §6.3 Limitations. The
performance experiments that follow support the quantitative argument ---
including the head-to-head vector benchmark against ObjectBox 4.x and SQLiteAI
sqlite-vector 0.9.95 in §5.8.

\subsection{Why Embed Valkey, and Why as a
Fork}\label{why-embed-valkey-and-why-as-a-fork}

Valkey is designed for Linux servers: long-lived daemon, TCP loopback, GNU libc
POSIX assumptions, shutdown via \texttt{exit()}. None of these survive a port to
a mobile app process. Dazzle applies three changes as a fork overlay: (a)
\textbf{portability} --- shims for Android Bionic libc and the Apple SDK; (b)
\textbf{lifecycle} --- the server cannot kill its host process; (c)
\textbf{in-process transport} --- TCP loopback is pure overhead when client and
server share the address space.

Upstream Valkey is not vendored; it is downloaded at a fixed tag and patched
with three diffs (\textasciitilde180 lines total) before compilation. Publishing
only the diffs preserves Valkey's BSD-3-Clause license; Dazzle's original code
(in \texttt{core/} and \texttt{sdk/}) is Apache-2.0.

\section{The Dazzle System}\label{the-dazzle-system}

\subsection{Architecture}\label{architecture}

Dazzle is organized in three layers:

\begin{itemize}
\tightlist
\item
  \textbf{\texttt{core/}} (Apache-2.0). In-process transport
  (\texttt{core/transport/dazzle\_transport.c}) + snapshot cache + worker pool
  (Android). Pure C, compiled with the patched Valkey.
\item
  \textbf{\texttt{versions/\textless{}ver\textgreater{}/patches/}} (build-time
  diffs). Three patches: Android (Bionic shims), iOS (Apple SDK adjustments + a
  custom \texttt{malloc\_zone\_t} to report isolated RSS), and a 14-line hook in
  \texttt{server.c} that inserts \texttt{dazzle\_direct\_init()} after
  \texttt{InitServerLast()}.
\item
  \textbf{\texttt{sdk/android/}, \texttt{sdk/ios/}} (Apache-2.0). JNI + Kotlin
  bridges (AAR) and C + Swift bridges (XCFramework).
\end{itemize}

\subsection{Transport Path}\label{transport-path}

Valkey's event loop runs on a dedicated \texttt{pthread}. The application thread
communicates via:

\begin{itemize}
\tightlist
\item
  \textbf{In-process pipe} --- the app thread \texttt{write()}s an 8-byte
  pointer to a \texttt{DirectRequest}, well below \texttt{PIPE\_BUF} (≥ 512
  bytes by POSIX, 4 096 on Linux), so the kernel guarantees atomic delivery
  without a write-side lock. The event loop wakes up on \texttt{read()},
  executes \texttt{call()} with a fake client, and signals the \texttt{condvar}
  with the RESP reply.
\item
  \textbf{Write-through snapshot cache} --- every mutating command crossing the
  pipe updates an in-process cache under rwlock before the caller is released.
  Subsequent reads acquire the rdlock and return values directly --- no
  event-loop wake-up. The pipe is synchronous, so by the time a read fires the
  cache reflects every prior write: \textbf{zero staleness by construction}.
\item
  \textbf{Worker pool (Android)} --- 2--4 additional threads absorb concurrent
  reads under per-slot striped rwlocks, when the environment flag
  \texttt{DAZZLE\_PARALLEL\_READS=1} is active.
\end{itemize}

\subsection{Post-EVAL Auto-Mirror}\label{post-eval-auto-mirror}

Lua scripts (\texttt{EVALSHA}) execute their internal writes through Valkey's
\texttt{call()}, which does not fire the snapshot-mirror hook. So that backends
writing fields inside a Lua script and reading them later don't pay the
pipe-HMGET cost on every read, Dazzle automatically hydrates the cache
post-EVAL: after \texttt{call()}, it iterates the script's declared KEYS,
re-reads each hash through the kvstore API, and upserts its fields into the
corresponding bucket. A backend that writes 8 fields inside a Lua EVALSHA sees
them reflected in the snapshot at the next read with no manual gymnastics from
the SDK.

\section{Experiment Design}\label{experiment-design}

\subsection{Workload}\label{workload}

Synthetic IoT monitoring agent: 200 readings with five fields (temperature,
humidity, pressure, battery, vibration), 11 injected anomalies, 10 inference
checkpoints evenly spaced. Two conditions: \textbf{Stateless} (model receives
only the last 20 readings) and \textbf{Augmented} (the 20 readings + a
materialized context block with aggregates over the full window).

\subsection{Backends}\label{backends}

Five embedded alternatives (SQLite, LMDB, ObjectBox, RocksDB, InMemory) plus a
Valkey-over-RESP-TCP reference (the unmodified Valkey 8 server with localhost
loopback, included as the ``transport overhead'' baseline) and seven Dazzle
variants exposing different primitive paths: basic, Lua, Pipeline, HFE, HLL,
\textbf{Precompute} (materialized aggregates) and \textbf{Vector} (HNSW vector
search). Twelve backends total per device, with \textbf{full bench-coverage
parity} across the storage-only and vector-bench sweeps in this revision (every
backend runs on both platforms; the ObjectBox 5.3 iOS port lands in §5.6, and
the Dazzle-Vector path runs on both platforms via the unified
\texttt{paper384\_scale} preset). This \emph{coverage} parity is distinct from
two limitations called out in §6.3 that the rest of the paper does not paper
over: (a) the Phase-2 transport optimisations (SPSC ring buffer,
\texttt{io\_uring}) are Linux-only --- iOS keeps the Phase-0 self-pipe path,
sufficient for the workload measured in §5.5; and (b) the end-to-end LLM
experiment of §5.9 runs on Moto G35 5G only.

\subsection{Devices}\label{devices}

\begin{itemize}
\tightlist
\item
  \textbf{Moto G35 5G} (Unisoc T760, 8×ARM @ 2.0/2.2 GHz, 4 GB LPDDR4X, Android
  14). Lower bound: \textasciitilde150 USD smartphone.
\item
  \textbf{iPhone 12 Pro} (Apple A14 Bionic, 6 cores, 6 GB, iOS 26.3). Runs the
  full twelve-backend sweep on physical hardware (no simulator), with one
  iOS-specific build adjustment: LMDB is opened with \texttt{MDB\_NOLOCK}
  because the iOS Documents sandbox blocks \texttt{fcntl(F\_SETLK)} on the LMDB
  lock file. Single-process agent --- lock coordination is moot --- so the flag
  does not affect the storage-only metrics.
\end{itemize}

\subsection{Metrics}\label{metrics}

Retrieval latency (mean, p50, p95, p99), ingest throughput, context block token
count, RAM delta, development complexity (LOC), model accuracy (anomaly recall
and factual synthesis accuracy on three verifiable global statistics).

\section{Evaluation}\label{evaluation}

\subsection{External Memory Enables
Synthesis}\label{external-memory-enables-synthesis}

\textbf{Table 2 --- Factual synthesis accuracy on global statistics (final
checkpoint).}

{\def\LTcaptype{none} % do not increment counter
\begin{longtable}[]{@{}
  >{\raggedright\arraybackslash}p{(\linewidth - 8\tabcolsep) * \real{0.2804}}
  >{\centering\arraybackslash}p{(\linewidth - 8\tabcolsep) * \real{0.2150}}
  >{\centering\arraybackslash}p{(\linewidth - 8\tabcolsep) * \real{0.1963}}
  >{\centering\arraybackslash}p{(\linewidth - 8\tabcolsep) * \real{0.2243}}
  >{\centering\arraybackslash}p{(\linewidth - 8\tabcolsep) * \real{0.0841}}@{}}
\toprule\noalign{}
\begin{minipage}[b]{\linewidth}\raggedright
Condition
\end{minipage} & \begin{minipage}[b]{\linewidth}\centering
Total readings
\end{minipage} & \begin{minipage}[b]{\linewidth}\centering
Anomaly count
\end{minipage} & \begin{minipage}[b]{\linewidth}\centering
Mean temperature
\end{minipage} & \begin{minipage}[b]{\linewidth}\centering
Score
\end{minipage} \\
\midrule\noalign{}
\endhead
\bottomrule\noalign{}
\endlastfoot
Stateless & -- (confabulates \textasciitilde50) & -- (confabulates 2--3) & --
(confabulates \textasciitilde21 °C) & \textbf{0/3} \\
Augmented (any backend) & Y (200) & Y (11) & Y (22.4 °C) & \textbf{3/3} \\
\end{longtable}
}

This is the central empirical finding of the paper. Without external memory the
model has no access to data outside its current context window (the last 20
readings); the ``global statistics'' it produces are confabulations. With a
materialized context block, it reads the exact values and reproduces them. The
specific backend is irrelevant for this finding --- what matters is that the
memory exists and that the context is precise.

The implication for agent design is direct: \textbf{the external memory layer is
mandatory infrastructure for any task that reasons about accumulated state}, not
an optional optimization. Backend choice then comes down to performance,
primitive surface, and ergonomics.

\subsection{Latency Envelope
Characterization}\label{latency-envelope-characterization}

The purpose of this section is to establish the latency envelope in which an
on-device agent backend has to operate, not to declare a winner on
microbenchmark numbers. We measure each backend in two regimes: a storage-only
sweep at N = 200 readings on both physical devices (Table 3), and a
concurrent-load reference at N = 20 000 (1 writer + 1 reader, Moto G35 5G)
carried over as a steady-state stress test. The two together cover the realistic
operating envelope of a long-running on-device agent.

\textbf{Methodological note for Tables 3--6.} SQLite-family rows
(\texttt{sqlite}, \texttt{sqlite-optimized}, \texttt{sqlite-precompute}) come
from an isolated mitigation sweep (Android: 3 rounds per backend, randomized
order; iOS: 3 rounds, dataset \texttt{dataset\_iot\_baseline}). Vector
SQLite-family rows (\texttt{sqlite-vec} and SQLiteAI variants) come from the
aligned vector harness at dim = 384, k = 10, p50 search latency over 100
queries. The iOS SQLiteAI path required vendoring the SQLite amalgamation with
\texttt{SQLITE\_ENABLE\_LOAD\_EXTENSION=1} into the benchmark target, because
Apple's system libsqlite3 ships with extension loading disabled --- this is
documented as part of the artifact (§8.1,
\texttt{experiment/backends/ios/sqlitevectorai/svai\_ios.c}) and was validated
end-to-end with the simulator-runnable XCTest suite before the device run.

\textbf{Table 3 --- Storage backends, N = 200 ingest/retrieval envelope on
physical devices} (per-reading ingest µs / average retrieval µs).

{\def\LTcaptype{none} % do not increment counter
\begin{longtable}[]{@{}
  >{\raggedright\arraybackslash}p{(\linewidth - 4\tabcolsep) * \real{0.2826}}
  >{\raggedleft\arraybackslash}p{(\linewidth - 4\tabcolsep) * \real{0.3478}}
  >{\raggedleft\arraybackslash}p{(\linewidth - 4\tabcolsep) * \real{0.3696}}@{}}
\toprule\noalign{}
\begin{minipage}[b]{\linewidth}\raggedright
Backend
\end{minipage} & \begin{minipage}[b]{\linewidth}\raggedleft
Moto G35 5G ingest / retrieval
\end{minipage} & \begin{minipage}[b]{\linewidth}\raggedleft
iPhone 12 Pro ingest / retrieval
\end{minipage} \\
\midrule\noalign{}
\endhead
\bottomrule\noalign{}
\endlastfoot
Dazzle (default) & 2 889 / 2 132 & 185 / 411 \\
Dazzle-Lua & 1 345 / 1 824 & 181 / 251 \\
Dazzle-Pipeline & 1 331 / 1 959 & 143 / 444 \\
Dazzle-HFE & 2 312 / 1 893 & 177 / 389 \\
Dazzle-HLL & 2 844 / 2 491 & 185 / 454 \\
\textbf{Dazzle-Precompute} & \textbf{1 404 / 34.7} & \textbf{267 / 129} \\
Valkey 8 (RESP-TCP) & 3 754 / 3 320 & 432 / 1 474 \\
SQLite (default) & 251.95 / 1 061.3 & 61.94 / 268.7 \\
SQLite-Optimized & 59.09 / 899.4 & 23.64 / 199.4 \\
SQLite-Precompute & 207.77 / 75.0 & 317.10 / 27.6 \\
LMDB & 172 / 641 & 100 / 346 \\
RocksDB & 294 / 922 & 169 / 424 \\
ObjectBox \(\ddagger\) & 1 839 / 778 & 4 840 / 638 \\
In-memory (struct) & 9.97 / 282 & 5.35 / 204 \\
\end{longtable}
}

\(\ddagger\) ObjectBox iOS ingest dropped from 24 780 µs/r to 4 840 µs/r (5.1×
speedup) once we wrapped each \texttt{ingest()} and
\texttt{storeCheckpointDecision()} body in
\texttt{store.runInTransaction\ \{\ ...\ \}}. The Kotlin port runs
\textasciitilde10 \texttt{box.put()} / \texttt{box.query()} calls per reading
and relies on ObjectBox-Java's implicit per-call write-tx batching to reach 1
839 µs/r; ObjectBox-Swift opens a fresh write transaction per call (lock +
journal append + fsync), so the same 10 calls without coalescing pay the lock
overhead 10×. The remaining \textasciitilde2.6× iPhone-vs-Moto gap is consistent
with the iOS allocator and FFI overhead observed elsewhere in the table (see the
Dazzle iOS footprint anomaly discussed in §6.3) --- not a bench-harness bug.

\textbf{Table 4 --- Vector backends, N = 200 ingest/retrieval envelope on
physical devices} (per-reading ingest µs / p50 search latency µs, dim = 384, k =
10).

{\def\LTcaptype{none} % do not increment counter
\begin{longtable}[]{@{}
  >{\raggedright\arraybackslash}p{(\linewidth - 4\tabcolsep) * \real{0.3103}}
  >{\raggedleft\arraybackslash}p{(\linewidth - 4\tabcolsep) * \real{0.3333}}
  >{\raggedleft\arraybackslash}p{(\linewidth - 4\tabcolsep) * \real{0.3563}}@{}}
\toprule\noalign{}
\begin{minipage}[b]{\linewidth}\raggedright
Backend
\end{minipage} & \begin{minipage}[b]{\linewidth}\raggedleft
Moto G35 5G ingest / search
\end{minipage} & \begin{minipage}[b]{\linewidth}\raggedleft
iPhone 12 Pro ingest / search
\end{minipage} \\
\midrule\noalign{}
\endhead
\bottomrule\noalign{}
\endlastfoot
\textbf{Dazzle-Vector (HNSW)} & \textbf{74 / 65} & \textbf{25.17 / 18} \\
sqlite-vec default & 23.62 / 916.0 & 27.64 / 196 \\
sqlite-vec optimized & 9.66 / 860.0 & 20.67 / 193 \\
sqlite-vec precompute & 9.35 / 816.7 & 19.72 / 181 \\
SQLiteAI default & 6.08 / 135.3 & 6.90 / 29 \\
SQLiteAI optimized & 5.98 / 135.0 & 6.82 / 29 \\
SQLiteAI precompute & 5.17 / 111.3 & 7.27 / 23 \\
\end{longtable}
}

Dazzle-Vector on Moto G35 5G now reports \textbf{74 µs/r ingest, 65 µs p50
search at recall 1.000 (HNSW, ef = 10)} at N = 200 --- the prior
\texttt{—\ /\ —} cell in this row was caused by a harness bug in
\texttt{SqliteBruteforceVector.create()} (the truth-source brute-force-SQLite
never wiped its database between runs because the file delete pointed at the
wrong directory, so the per-config top-k truth was computed against
\textasciitilde20 000 stale docs accumulated across previous bench launches;
only the last config in any sweep fortuitously had recall recover above the 0.95
floor). Both fixes land in this revision: the truth-source delete now resolves
to \texttt{context.getDatabasePath(...)} and \texttt{VectorBenchmark} issues
\texttt{FLUSHALL} between configs and between Dazzle engine variants so each
measurement is per-engine and per-config in isolation. The fix landed in commit
\texttt{1e3d5f5} of \texttt{experiment/backends/android/core/}
(\texttt{SqliteBruteforceVector.kt} + \texttt{VectorBenchmark.kt}), and the
post-fix benchmark run is committed at
\texttt{research/benchmarks/results/Moto\_G35\_5G/vector/vecbench\_moto\_g35\_5G\_1777369156656.json}
(SHA-256
\texttt{5a64d3692da166d96306d456697e43bb89c27c07b515e253346c5b33bc5c9b5b}).
Tables 4, 5, 11 and the §5.8.4 / companion-report N-sweep all read from this
single JSON, so the cell-by-cell numbers and the recall narrative come from one
device run, not from a mix of pre-fix and post-fix snapshots.

Observation 1: Dazzle-Precompute sets the floor reachable when materialized
aggregates are exposed as a native primitive. Retrieval averages 34.7 µs at N =
200 on Moto G35 5G and 47 µs at N = 20 000 under concurrent load --- flat in N
because the materialized context block is rebuilt on the write path and read
with a constant-time field copy. On iPhone 12 Pro retrieval lands at 129 µs.
These numbers represent the lower bound a backend can reach when the agent's
read pattern is compiled into the schema, not what any backend ``could'' achieve
in principle.

Observation 2: honest concession, now reflected in the main table with all nine
SQLite-family variants (default/optimized/precompute across storage SQLite,
sqlite-vec, and SQLiteAI). For the storage path, Android retrieval moves from 1
061.3 us (\texttt{sqlite} default) to 899.4 us (\texttt{sqlite-optimized}) and
to 75.0 us (\texttt{sqlite-precompute}). This closes most of the storage-latency
gap and reinforces the paper's framing: the key differentiator is not an
intrinsic SQLite algorithm limit, but how much materialization ceremony and
maintenance code the application must own.

Observation 3: the observation that actually matters. Even the highest retrieval
latency in the envelope (SQLite at 3 374 µs under concurrent load on Moto G35
5G) is roughly 0.0001 \% of a single Gemma 4 E2B \citep{gemma4} inference turn
(1-3 seconds). The user-perceived difference between 50 µs and 3 374 µs at the
agent layer is zero: both are dwarfed by the inference call that sits
immediately after the retrieval. Microsecond comparisons are therefore not the
right axis on which to choose a backend. The dominant criteria are the primitive
surface (Table 1, §2.2) and the development complexity required to deliver a
working state layer (§5.6).

\subsubsection{Why we still report exact microsecond
numbers}\label{why-we-still-report-exact-microsecond-numbers}

A reviewer might reasonably ask: if every retrieval is below the inference noise
floor and microsecond comparisons are not the choice axis, why does §5 spend
twelve tables enumerating microseconds? Two reasons, both independent of the
agent thesis:

\textbf{(i) Non-LLM consumers may share the storage layer.} Although this paper
does not benchmark them, the same Dazzle code path is structurally compatible
with on-device analytics dashboards, sensor pre-aggregation pipelines, and
sync-conflict resolution layers --- workloads where the retrieval primitives
(materialised aggregates, bounded streams, HyperLogLog, HNSW) are the same but
where retrieval latency is \emph{not} dominated by an LLM forward pass. We
report the absolute microsecond numbers so a reader evaluating those workloads
can pull the relevant row out of Table 3 / Table 5 / Table 11 directly without
re-running the bench. We make no competitive claim about Dazzle's position in
those workloads at this time.

\textbf{(ii) Absolute numbers are a sanity check on the architecture.} A backend
that drifts to 100 ms+ at N = 200 (an order of magnitude above what we report
here) would suggest a regression worth investigating regardless of the LLM noise
floor --- the Write-Heavy / Read-Light pattern is supposed to keep retrieval
bounded; if a future revision breaks that, the exact microsecond column in T3 is
where the regression shows up. Same logic for the §5.4 ablation and the §5.7
performance-evolution table in the companion report: those numbers are how we
know the architectural pattern is intact across revisions, not how we recommend
choosing a backend.

This section therefore establishes that all measured backends sit within the
acceptable envelope for an on-device LLM agent. The remainder of the evaluation
reads on two tracks: \textbf{(a)} for an agent-loop reader, primitive surface
(Table 1) and development complexity (Table 9) drive backend selection;
\textbf{(b)} for a non-LLM reader, the per-engine microsecond numbers are the
direct input. Both readings come from the same data; the difference is which
axis the reader should treat as the headline.

\subsection{Constant-N Retrieval and Storage
Footprint}\label{constant-n-retrieval-and-storage-footprint}

\textbf{Table 5 --- Retrieval latency average across dataset sizes, Moto G35 5G}
(µs). Sweep at N \(\in\) \{200, 1 000, 5 000, 20 000\}. Each cell in Table 5 and
Table 5b is the per-run arithmetic mean of the retrieval-sample window
(\texttt{retrieval\_avg\_us} in each
\texttt{scale\_\textless{}backend\textgreater{}\_*.json}); for backends with
multiple sweep runs in the repo (typically 3--4), the cell is the median of
those per-run means. The corresponding \texttt{retrieval\_p50\_us} column is
also persisted in every JSON for any reader who prefers the median over the
mean, but the table values throughout this section are means.

{\def\LTcaptype{none} % do not increment counter
\begin{longtable}[]{@{}lrrrr@{}}
\toprule\noalign{}
Backend & N = 200 & N = 1 k & N = 5 k & N = 20 k \\
\midrule\noalign{}
\endhead
\bottomrule\noalign{}
\endlastfoot
\textbf{Dazzle-Precompute} & \textbf{37} & \textbf{29} & \textbf{29} &
\textbf{34} \\
\textbf{Dazzle-Incremental} & 325 & 163 & 155 & 146 \\
InMemory & 663 & 314 & 273 & 317 \\
SQLite (default) & 1 061.3 & 739.1 & 717.4 & 721.6 \\
SQLite-Optimized & 899.4 & 508.1 & 514.5 & 505.4 \\
SQLite-Precompute & 75.0 & 59.8 & 62.1 & 60.6 \\
sqlite-vec default & 916.0 & 1 445.7 & 7 271.0 & 27 850.3 \\
sqlite-vec optimized & 860.0 & 1 434.7 & 7 145.3 & 28 574.7 \\
sqlite-vec precompute & 816.7 & 1 430.7 & 7 421.3 & 26 505.3 \\
SQLiteAI default & 135.3 & 303.0 & 1 475.0 & 9 812.3 \\
SQLiteAI optimized & 135.0 & 302.3 & 1 591.7 & 9 569.7 \\
SQLiteAI precompute & 111.3 & 233.3 & 889.0 & 3 071.7 \\
LMDB & 1 130 & 841 & 1 398 & 1 986 \\
\end{longtable}
}

\textbf{Table 5b --- Retrieval latency average across dataset sizes, iPhone 12
Pro} (µs). Same metric (arithmetic mean of per-sample retrieval latency,
\texttt{retrieval\_avg\_us}) and same sweep grid as Table 5.

{\def\LTcaptype{none} % do not increment counter
\begin{longtable}[]{@{}lrrrr@{}}
\toprule\noalign{}
Backend & N = 200 & N = 1 k & N = 5 k & N = 20 k \\
\midrule\noalign{}
\endhead
\bottomrule\noalign{}
\endlastfoot
\textbf{Dazzle-Precompute} & \textbf{25.4} & \textbf{17.8} & \textbf{17.8} &
\textbf{17.8} \\
\textbf{Dazzle-Incremental} & 139.2 & 63.7 & 63.3 & 63.2 \\
InMemory & 148.9 & 149.3 & 106.4 & 55.8 \\
SQLite (default) & 73.4 & 47.5 & 47.7 & 46.7 \\
SQLite-Optimized & 75.9 & 38.2 & 32.9 & 33.1 \\
SQLite-Precompute & 6.88 & 3.75 & 3.79 & 3.83 \\
sqlite-vec default & 196 & 548 & 3 578 & 15 038 \\
sqlite-vec optimized & 193 & 547 & 3 535 & 15 200 \\
sqlite-vec precompute & 181 & 548 & 3 535 & 15 145 \\
SQLiteAI default & 29 & 96 & 441 & 2 850 \\
SQLiteAI optimized & 29 & 96 & 440 & 2 842 \\
SQLiteAI precompute & 23 & 79 & 355 & 1 407 \\
LMDB & 167.4 & 67.2 & 46.9 & 85.8 \\
\end{longtable}
}

Tables 5 and 5b should be read as an implementation-pattern result, not as a
backend horse race. The relevant property is \textbf{constant-in-N retrieval}
when the system serves materialized aggregates instead of rebuilding context
from raw history at query time. The Dazzle-Precompute row in Table 5 is the
median of the per-run means from four sweep runs (three archived in
\texttt{research/benchmarks/results/motorola\_moto\_g35\_5G/dazzle-precompute/scale\_dazzle-precompute\_*.json}
and one re-run dated 2026-04-28 archived in
\texttt{research/benchmarks/results/Moto\_G35\_5G/scale\_dazzle\_precompute\_rerun/scale\_dazzle-precompute\_1777401499448.json});
v1 of this paper rendered the N = 5 k cell as \textbf{77 µs} --- a single
mis-transcribed value (the median of per-run means across all four runs is 28.71
µs, and the corresponding median of per-run p50s is 27.27 µs); the value
reported in the row above is the corrected median-of-means rounded to integer
microseconds.

Dazzle-Precompute and Dazzle-Incremental provide this pattern natively:
aggregation work is paid on write, and retrieval reads bounded materialized
state. SQLite can also realize near-constant retrieval when the same pattern is
implemented explicitly (for example, through materialized scalar rows and
indexed B-tree lookups). Constant-in-N behavior is therefore a property of the
materialized-aggregate approach, not a proprietary property of one engine.

LMDB and InMemory, in the default forms benchmarked here, follow a
recompute-on-read path for this workload and therefore do not provide the same
architectural constant-in-N guarantee. LMDB shows clear growth with N in Table
5; InMemory shows lower absolute numbers at this scale but still relies on
read-time reconstruction rather than a persisted materialized-aggregate
primitive.

The strongest result in this section is about prompt behavior, not raw speed
ranking: the materialized context block keeps a \textbf{bounded token budget} of
about \textbf{58--62 tokens} independent of N. Unmaterialized paths that return
broader raw fields typically produce much larger payloads (often 200+ tokens),
which directly affects prompt budget, prefill cost, and long-run interaction
stability.

\subsubsection{Storage footprint after 200
ingests}\label{storage-footprint-after-200-ingests}

The same physical-device sweep also produces a clean separation on on-disk
footprint, measured per backend after the 200-reading run finishes
(\texttt{stat.st\_blocks\ ×\ 512}, matches \texttt{du\ -k}).

\textbf{Table 6 --- On-disk footprint after 200 ingests, by device} (KB).

{\def\LTcaptype{none} % do not increment counter
\begin{longtable}[]{@{}lrr@{}}
\toprule\noalign{}
Backend & Moto G35 5G (KB) & iPhone 12 Pro (KB) \\
\midrule\noalign{}
\endhead
\bottomrule\noalign{}
\endlastfoot
\textbf{Dazzle (default)} & \textbf{6.7} & \textbf{153.3} \\
Dazzle-Lua & 6.6 & 153.3 \\
Dazzle-Pipeline & 6.7 & 153.3 \\
Dazzle-HFE & 6.6 & 153.3 \\
Dazzle-HLL & 7.0 & 177.6 \\
Dazzle-Precompute & 7.7 & 153.5 \\
Valkey 8 (RESP-TCP) & 7.9 & 177.4 \\
LMDB & 72.0 & 176.0 \\
ObjectBox & 96.0 & 268.0 \\
RocksDB & 4 515.8 & 104.0 \\
SQLite & 374.2 & 3 635.2 \\
SQLite-Optimized & 494.6 & 3 391.5 \\
SQLite-Precompute & 495.6 & 176.6 \\
sqlite-vec default & 4.1 & 4.0 \\
sqlite-vec optimized & 4.1 & 4.0 \\
sqlite-vec precompute & 4.1 & 4.0 \\
SQLiteAI default & 12.3 & 12.0 \\
SQLiteAI optimized & 12.3 & 12.0 \\
SQLiteAI precompute & 12.3 & 12.0 \\
\end{longtable}
}

The headline ratio is \textbf{Dazzle vs SQLite}:

\begin{itemize}
\tightlist
\item
  Moto G35 5G: 6.7 KB vs 374.2 KB --- \textbf{56× smaller}
\item
  iPhone 12 Pro: 153.3 KB vs 3 635.2 KB --- \textbf{24× smaller}
\end{itemize}

The size advantage is structural, not incidental. Dazzle stores a flat hash with
five aggregate fields plus a bounded suffix of the most recent 200 readings (cap
= 200). SQLite, by contrast, allocates per-row B-tree pages plus per-table index
pages, so the on-disk footprint grows roughly with N even after the container
metadata is amortized. RocksDB on Android shows the same problem from a
different angle (4.4 MB after 200 reads --- SST file headers alone exceed the
entire Dazzle payload). LMDB lands in the middle: a single page-aligned
\texttt{data.mdb} file with no indices, giving 72 KB on Moto and 176 KB on
iPhone.

Footprint is the binding constraint on IoT and budget mobile hardware where
flash budget and DRAM bandwidth are both scarce. A \textbf{24-56× smaller}
payload is the difference between Dazzle's data fitting in a handful of L2 cache
lines per retrieval versus sweeping multiple SQLite pages per query --- which
compounds back into the retrieval-latency picture of §5.2.

\textbf{Two platform-specific footprint outliers warrant honest disclosure.}
First, Dazzle on iOS reports 153 KB against 6.7 KB on Android for the
\emph{same} in-process Valkey-fork running the \emph{same} 200-reading workload
--- a 23× cross-platform skew. The metric (\texttt{used\_memory\_dataset}) and
the input data are identical; the gap is in the Valkey allocator behavior under
Apple's libsystem \texttt{malloc} versus the \texttt{jemalloc} arena Android
uses. Even at the larger iOS figure, Dazzle still beats SQLite by 24× on the
same device, so the headline ratio holds; we flag the iOS allocator footprint as
a future-investigation item (§6.3). Second, RocksDB shows a 42× skew in the
\emph{opposite} direction (4.41 MB on Android, 104 KB on iOS) driven by Android
defaults emitting a multi-MB WAL plus initial SST file pre- allocation, while
the iOS port uses RocksDB's ``minimal'' default slice. Both anomalies sharpen
the same point: native database footprint on mobile is a configuration-sensitive
metric, and any single-number comparison should be read as ``best-effort
default'' rather than ``tuned floor''.

\subsection{Factorial 2³ Ablation: Which Layer
Contributes?}\label{factorial-2uxb3-ablation-which-layer-contributes}

The Dazzle stack composes three independent optimizations: snapshot cache
(on/off), bucket dispatch by hash (1 bucket vs 16), parallel worker pool
(off/on). We measure six valid combinations on the same load --- K = 8 agents,
80/20 read/write mix, 15 s per cell.

\textbf{Table 7 --- Throughput at K = 8, Moto G35 5G} (ops/s, 80/20 read/ write
mix, 15 s per cell).

{\def\LTcaptype{none} % do not increment counter
\begin{longtable}[]{@{}
  >{\raggedright\arraybackslash}p{(\linewidth - 12\tabcolsep) * \real{0.2651}}
  >{\raggedleft\arraybackslash}p{(\linewidth - 12\tabcolsep) * \real{0.1205}}
  >{\raggedleft\arraybackslash}p{(\linewidth - 12\tabcolsep) * \real{0.1084}}
  >{\raggedleft\arraybackslash}p{(\linewidth - 12\tabcolsep) * \real{0.1446}}
  >{\raggedleft\arraybackslash}p{(\linewidth - 12\tabcolsep) * \real{0.1205}}
  >{\raggedleft\arraybackslash}p{(\linewidth - 12\tabcolsep) * \real{0.1205}}
  >{\raggedleft\arraybackslash}p{(\linewidth - 12\tabcolsep) * \real{0.1205}}@{}}
\toprule\noalign{}
\begin{minipage}[b]{\linewidth}\raggedright
Backend
\end{minipage} & \begin{minipage}[b]{\linewidth}\raggedleft
baseline
\end{minipage} & \begin{minipage}[b]{\linewidth}\raggedleft
workers
\end{minipage} & \begin{minipage}[b]{\linewidth}\raggedleft
snap-ser
\end{minipage} & \begin{minipage}[b]{\linewidth}\raggedleft
snap-par
\end{minipage} & \begin{minipage}[b]{\linewidth}\raggedleft
hash-ser
\end{minipage} & \begin{minipage}[b]{\linewidth}\raggedleft
hash-par
\end{minipage} \\
\midrule\noalign{}
\endhead
\bottomrule\noalign{}
\endlastfoot
\texttt{dazzle-incremental} & 12 106 & 9 604 & \textbf{28 897} & 28 154 & 27 926
& 28 304 \\
\texttt{dazzle-precompute} & 19 676 & 10 988 & \textbf{38 830} & 38 754 & 38 367
& 38 156 \\
\end{longtable}
}

Three quantified observations:

\begin{enumerate}
\def\labelenumi{\arabic{enumi}.}
\tightlist
\item
  \textbf{Snapshot cache alone is the throughput layer.}
  \texttt{snap-linear-serial} (snapshot on, 1 bucket, no workers) is the peak in
  both rows: \textbf{38 830 ops/s} for precompute and \textbf{28 897 ops/s} for
  incremental. Layering the 16-bucket hash index and the parallel worker pool on
  top of the cache (the \texttt{hash-par} ``full-stack'' column) actually
  delivers \textbf{38 156 / 28 304 ops/s} --- i.e.~snap-ser is 1.8 \% above
  hash-par for precompute and 2.1 \% above for incremental. The full stack
  loses, not gains, against snap-ser; the deltas are within run-to-run variance,
  so the practical reading is ``the snapshot cache by itself is sufficient and
  the other two layers add nothing measurable on this workload.''
\item
  \textbf{Workers alone regress performance.} Without a cache to dispatch
  against, the worker pool is pure queueing overhead: 44 \% slower for
  precompute, 21 \% slower for incremental.
\item
  \textbf{Hash-index is future-proofing.} The 4--5-key keyspace of the benchmark
  fits trivially in a single bucket; the O(1) dispatch becomes visible only when
  the keyspace grows to dozens of keys or when field names share long prefixes.
\end{enumerate}

The clean narrative: \textbf{the snapshot cache is the only layer that moves
throughput on this workload; the parallel worker pool and the hash-index
dispatch are no-ops here and only earn their cost on workloads with a larger or
more contended keyspace}.

\subsection{Cross-Platform Validation: iPhone 12
Pro}\label{cross-platform-validation-iphone-12-pro}

Same Dazzle stack run on a physical iPhone 12 Pro (iOS 26.3, A14 Bionic) and a
physical Moto G35 5G (Android 14, Unisoc T760), 200 readings, single writer +
100 retrieval samples (same storage-only protocol as §5.2).

\textbf{Table 8 --- Cross-platform Dazzle summary, Moto G35 5G vs iPhone 12 Pro}
(N = 200, single writer + 100 retrieval samples).

{\def\LTcaptype{none} % do not increment counter
\begin{longtable}[]{@{}lrr@{}}
\toprule\noalign{}
Metric & Moto G35 5G & iPhone 12 Pro \\
\midrule\noalign{}
\endhead
\bottomrule\noalign{}
\endlastfoot
Dazzle-Precompute ingest (µs/r) & 1 404 & 267 \\
Dazzle-Precompute retrieval avg & \textbf{34.7 µs} & \textbf{129 µs} \\
Dazzle-Precompute retrieval P50 & 34.3 µs & 139 µs \\
Dazzle-Precompute retrieval P95 & 42.2 µs & 189 µs \\
Dazzle-default retrieval avg & 2 132 µs & 411 µs \\
Retrieval gain (precompute/default) & 61× & 3.2× \\
Footprint (Dazzle default) & 6.7 KB & 153.3 KB \\
Footprint vs SQLite & 56× smaller & 24× smaller \\
\end{longtable}
}

Three observations on cross-platform behavior:

\begin{enumerate}
\def\labelenumi{\arabic{enumi}.}
\item
  \textbf{Architectural shape generalizes.} Dazzle-Precompute beats the default
  Dazzle path on both platforms --- the Write-Heavy / Read-Light trade-off holds
  in both directions of the SoC + OS pair. The snapshot cache pays off on every
  read regardless of where it runs.
\item
  \textbf{Absolute retrieval numbers are platform-sensitive.} Moto G35 5G lands
  at 34 µs P50; iPhone 12 Pro at 139 µs P50. The 4× gap is driven by the
  in-process Valkey-fork's RESP path on iOS (tokio runtime + Rust→Swift bridge
  per retrieval), not by the snapshot cache itself. The Android path uses a
  thinner JNI bridge with fewer hops per call.
\item
  \textbf{Ingest direction reverses.} Moto needs 1 404 µs per reading; iPhone
  needs 267 µs (5.3× cheaper on iPhone). The A14 Bionic's faster cores plus
  LPDDR4X bandwidth swallow the per-reading \texttt{XADD} +
  \texttt{HINCRBYFLOAT} + \texttt{HSET} sequence faster than the Unisoc T760
  budget tier.
\end{enumerate}

The bottom line: the \textbf{architectural contribution generalizes} across
hardware and OS (cost shifts from read to write; snapshot cache absorbs N;
footprint stays at single-engine scale), while \textbf{absolute numbers vary by
SoC and runtime stack}. Even on the slower iPhone retrieval path, P95 stays
under 200 µs --- three to four orders of magnitude under the on-device LLM
forward pass (\(>10^{6}\) µs on either device), which remains the only term that
matters at the agent loop level.

\subsection{Development Complexity}\label{development-complexity}

\textbf{Table 9 --- Lines of code per backend, broken out by platform} (SLOC:
blanks and pure-comment lines excluded; Kotlin counts use the Android
\texttt{*ContextManager.kt} file at \texttt{experiment/backends/android/} and
Swift counts use the iOS \texttt{*ContextManager.swift} file at
\texttt{experiment/backends/ios/}. Lower is better; every row implements the
same agent state-layer surface).

{\def\LTcaptype{none} % do not increment counter
\begin{longtable}[]{@{}
  >{\raggedright\arraybackslash}p{(\linewidth - 8\tabcolsep) * \real{0.1389}}
  >{\raggedleft\arraybackslash}p{(\linewidth - 8\tabcolsep) * \real{0.2130}}
  >{\raggedleft\arraybackslash}p{(\linewidth - 8\tabcolsep) * \real{0.1667}}
  >{\raggedleft\arraybackslash}p{(\linewidth - 8\tabcolsep) * \real{0.1111}}
  >{\raggedright\arraybackslash}p{(\linewidth - 8\tabcolsep) * \real{0.3704}}@{}}
\toprule\noalign{}
\begin{minipage}[b]{\linewidth}\raggedright
Backend
\end{minipage} & \begin{minipage}[b]{\linewidth}\raggedleft
Android (Kotlin) SLOC
\end{minipage} & \begin{minipage}[b]{\linewidth}\raggedleft
iOS (Swift) SLOC
\end{minipage} & \begin{minipage}[b]{\linewidth}\raggedleft
Sum (both)
\end{minipage} & \begin{minipage}[b]{\linewidth}\raggedright
API style
\end{minipage} \\
\midrule\noalign{}
\endhead
\bottomrule\noalign{}
\endlastfoot
\textbf{Dazzle} & \textbf{175} & \textbf{186} & \textbf{361} & Typed primitive
commands \\
InMemory & 174 & 172 & 346 & Direct collection manipulation \\
LMDB & 200 & 227 & 427 & Byte-buffer API \\
RocksDB & 203 & 227 & 430 & Byte-buffer API with options \\
ObjectBox & 233 & 245 & 478 & Schema annotations + DSL \\
SQLite & 268 & 292 & 560 & SQL + cursor mapping \\
\end{longtable}
}

The two-platform breakdown matters because the Android port and the iOS port of
a backend track each other within \textasciitilde30 SLOC for every row except
SQLite (+24 on iOS, mostly C-API result-row marshalling that Android's
\texttt{Cursor} API hides). The relative ordering is stable across platforms:
Dazzle is the lowest non-InMemory backend on both, SQLite is the highest, and
the gap between Dazzle and SQLite is \textasciitilde93 SLOC on Android and
\textasciitilde106 SLOC on iOS --- slightly wider on iOS.

Dazzle requires the fewest lines because its typed API (\texttt{XADD},
\texttt{HMGET}, \texttt{PFADD}) directly expresses the agent-domain operations
without schema translation or cursor parsing.

\subsection{Performance Evolution
(summary)}\label{performance-evolution-summary}

The current stack is the cumulative result of five in-process transport
redesigns. The single largest contribution is the \textbf{post-EVAL auto-mirror
(+189 \% on incremental retrieval throughput} at K = 8, Moto G35 5G, 80/20
read-write mix), which hydrates the snapshot cache from Lua-internal writes
without manual ceremony. Adding the bucketed snapshot index, the parallel worker
pool, and the inline-HSET in the Lua script contributes a further
\textasciitilde14 \% on precompute. The full version-by-version table (baseline
9 684 ops/s → final 28 057 ops/s on incremental, 34 356 → 38 156 on precompute)
lives in §1 of the companion engineering report
(\texttt{research/paper/companion\_engineering\_report.md}). Each optimization
was verified as an independent contribution; the final stack has lived in
\texttt{main} since April 2026.

\subsection{Vector Search Primitive: Envelope and Operating
Points}\label{vector-search-primitive-envelope-and-operating-points}

The headline result of this section is an envelope claim, not a speed
competition claim: \textbf{Dazzle's HNSW vector search delivers sub-millisecond
retrieval at recall \(\geq\) 0.95 across the standard mobile RAG operating range
(dim \(\leq\) 384, N \(\leq\) 10 000)}. That places Dazzle in the same
algorithmic class --- HNSW \citep{hnsw}, with expected \(O(\log N)\) search
complexity --- as ObjectBox 4.x, a commercial mobile-first peer with native HNSW
support.

SQLiteAI sqlite-vector v0.9.95, the production-shipping commercial SQLite
extension as of Q1 2026, exposes a SIMD-accelerated quantized linear scan
(\texttt{vector\_quantize\_scan}) as its only optimized query path; HNSW is not
user-accessible in this version. Including SQLiteAI is informative as a
reference for what mobile developers in the SQLite ecosystem encounter today,
but this comparison spans two algorithm classes (O(log N) vs O(N)); the
resulting gap reflects that asymmetry, not engineering quality of either
implementation.

The harness contract for every recall floor / operating-point cell in this
section is described, audited, and post-mortemed in Appendix B
(\texttt{research/paper/appendix\_b\_harness\_postmortem.md}). All cells in
Tables 4, 11 and 12 below come from the post-fix run archived at
\texttt{research/benchmarks/results/Moto\_G35\_5G/vector/vecbench\_moto\_g35\_5G\_1777369156656.json}
(commit \texttt{1e3d5f5}, SHA-256 cited in §5.2).

\subsubsection{5.8.1 Setup --- two harness
presets}\label{setup-two-harness-presets}

The vector evaluation actually runs the same harness
(\texttt{VectorBenchmark.kt} on Android, \texttt{VectorBenchmark.swift} on iOS)
against two different config grids, and the rest of §5.8 cites both depending on
the question being answered:

{\def\LTcaptype{none} % do not increment counter
\begin{longtable}[]{@{}
  >{\raggedright\arraybackslash}p{(\linewidth - 6\tabcolsep) * \real{0.3947}}
  >{\raggedleft\arraybackslash}p{(\linewidth - 6\tabcolsep) * \real{0.1184}}
  >{\raggedright\arraybackslash}p{(\linewidth - 6\tabcolsep) * \real{0.2632}}
  >{\raggedright\arraybackslash}p{(\linewidth - 6\tabcolsep) * \real{0.2237}}@{}}
\toprule\noalign{}
\begin{minipage}[b]{\linewidth}\raggedright
Preset
\end{minipage} & \begin{minipage}[b]{\linewidth}\raggedleft
Configs
\end{minipage} & \begin{minipage}[b]{\linewidth}\raggedright
Used in
\end{minipage} & \begin{minipage}[b]{\linewidth}\raggedright
What it answers
\end{minipage} \\
\midrule\noalign{}
\endhead
\bottomrule\noalign{}
\endlastfoot
\texttt{DEFAULT\_CONFIGS} & 9 (3 dim \(\times\) 3 N) & §5.8.3 envelope claim &
``Does Dazzle stay sub-millisecond across the standard mobile RAG range (dim
\(\leq\) 384, N \(\leq\) 10 000)?'' \\
\texttt{vector-bench-paper384-scale} (Android) / \texttt{paper384\_scale} (iOS)
& 4 (1 dim \(\times\) 4 N) & T11 + T12 + T13 + T14 & ``What is the operating
point at the paper's headline scale, and how does it scale to N = 20 000?'' \\
\end{longtable}
}

The 9-config grid sweeps \textbf{\{16, 128, 384\}} dim \(\times\) \textbf{\{500,
2 000, 10 000\}} N to characterise the shape of the envelope. The 4-config grid
pins dim = 384 and sweeps N \(\in\) \textbf{\{200, 1 000, 5 000, 20 000\}} so
the operating point reported in T11 (N = 20 000) is on the same grid as the
SQLite-family N-sweep in T12--T14. Both grids run k = 10 nearest neighbours,
cosine distance, recall floor recall@10 \(\geq\) 0.95, p50 / p95 / p99 over 100
queries per cell in warm steady-state (cold cache discarded). Where the text
below refers to \textbf{``the 9-config grid''}, it cites
\texttt{DEFAULT\_CONFIGS}; the operating-point table cites the paper384\_scale
run.

Configurations:

\begin{itemize}
\tightlist
\item
  \textbf{Dazzle HNSW (float32)}: M=32, efConstruction=400, efRuntime varied in
  \{10, 50, 100, 200\}.
\item
  \textbf{Dazzle SQ8 (int8 + NEON SDOT)}: same HNSW build params, plus scalar
  quantization on the storage path with NEON dot-product on the comparison path.
\item
  \textbf{ObjectBox 4.x}: HNSW via the \texttt{@HnswIndex} annotation,
  configured with the same M=32, efConstruction=400, cosine distance, float32
  vectors.
\item
  \textbf{SQLiteAI 0.9.95}: \texttt{vector\_init} with FLOAT32 + cosine,
  \texttt{vector\_quantize(max\_memory=50MB)} to build the int8 snapshot,
  \texttt{vector\_quantize\_preload} to warm the cache, queries through
  \texttt{vector\_quantize\_scan} (linear scan over quantized vectors).
\end{itemize}

To make the SQLite comparison auditable beyond a single extension path, this
revision also includes a dedicated SQLite-family sweep (§5.8.4) with explicit
variants for \texttt{sqlite-vec} and SQLiteAI: \texttt{default},
\texttt{optimized}, and \texttt{precompute} operating points.

\subsubsection{5.8.2 Operating-Point Table}\label{operating-point-table}

\textbf{Table 11 --- Vector engine operating-point comparison at dim = 384, N =
20 000, k = 10, recall@10 floor \(\geq\) 0.95} (p50 search latency in µs, ingest
total in ms, DB size on disk for SQLite-family engines,
\texttt{INFO\ memory\ →\ used\_memory\_dataset} for the in-memory Dazzle/Valkey
rows).

{\def\LTcaptype{none} % do not increment counter
\begin{longtable}[]{@{}
  >{\raggedright\arraybackslash}p{(\linewidth - 12\tabcolsep) * \real{0.2131}}
  >{\raggedright\arraybackslash}p{(\linewidth - 12\tabcolsep) * \real{0.1885}}
  >{\raggedright\arraybackslash}p{(\linewidth - 12\tabcolsep) * \real{0.0902}}
  >{\raggedleft\arraybackslash}p{(\linewidth - 12\tabcolsep) * \real{0.0902}}
  >{\raggedleft\arraybackslash}p{(\linewidth - 12\tabcolsep) * \real{0.1148}}
  >{\raggedleft\arraybackslash}p{(\linewidth - 12\tabcolsep) * \real{0.1066}}
  >{\raggedleft\arraybackslash}p{(\linewidth - 12\tabcolsep) * \real{0.1967}}@{}}
\toprule\noalign{}
\begin{minipage}[b]{\linewidth}\raggedright
Engine / Variant
\end{minipage} & \begin{minipage}[b]{\linewidth}\raggedright
Algorithm
\end{minipage} & \begin{minipage}[b]{\linewidth}\raggedright
Precision
\end{minipage} & \begin{minipage}[b]{\linewidth}\raggedleft
Recall@10
\end{minipage} & \begin{minipage}[b]{\linewidth}\raggedleft
p50 search ‡
\end{minipage} & \begin{minipage}[b]{\linewidth}\raggedleft
Ingest (ms)
\end{minipage} & \begin{minipage}[b]{\linewidth}\raggedleft
RAM / DB size
\end{minipage} \\
\midrule\noalign{}
\endhead
\bottomrule\noalign{}
\endlastfoot
\textbf{Dazzle SQ8} & HNSW + int8 SDOT & int8 & 0.959 & \textbf{208 µs} {[}203,
212{]} & 16 614 & 9.77 MB \(\dagger\) \\
Dazzle SQ8+Rerank & HNSW + int8 + fp32 & int8/fp32 & 0.982 & 330 µs & 14 966 &
39.06 MB \(\dagger\) \\
Dazzle F16 & HNSW + fp16 & fp16 & 0.984 & 297 µs & 26 082 & 17.09 MB
\(\dagger\) \\
Dazzle HNSW & HNSW & float32 & 0.952 & 343 µs & 43 028 & 31.74 MB \(\dagger\) \\
ObjectBox 4.x & HNSW & float32 & 0.994 & 853 µs {[}853, 1 078{]} & 387 405 & not
reported by runner \\
sqlite\_plain & linear scan & float32 & 1.000 & 707 302 µs & 1 614 & 79.45 MB \\
sqlite\_vec\_default & linear scan & float32 & 1.000 & 27 391 µs & 1 316 & 31.00
MB \\
sqlite\_vec\_optimized & linear scan & float32 & 1.000 & 28 575 µs & 869 & 31.00
MB \\
sqlite\_vec\_precompute & linear scan & float32 & 1.000 & 26 505 µs & 870 &
31.00 MB \\
SQLiteAI default \(*\) & quantized linear scan & int8 & 0.987 & 3 087 µs & 723 &
46.65 MB \\
SQLiteAI optimized \(*\) & quantized linear scan & int8 & 0.987 & 2 842 µs & 685
& 46.65 MB \\
SQLiteAI precompute \(*\) & quantized linear scan & int8 & 0.987 & 1 407 µs &
708 & 46.65 MB \\
\end{longtable}
}

‡ The bracketed \texttt{{[}lo,\ hi{]}} interval next to the two same-class
headline cells (Dazzle SQ8, ObjectBox 4.x --- both HNSW) is a \textbf{cross-run
percentile bootstrap 95 \% CI on the p50 statistic}
\citep{efronbootstrap, davisonhinkley} computed over the four independent
post-fix Moto G35 5G bench runs archived in
\texttt{research/benchmarks/results/Moto\_G35\_5G/vector/} (timestamps
2026-04-28T05:05Z / 05:21Z / 09:24Z / 09:39Z; B = 10 000, seed = 42; script
\texttt{research/scripts/bootstrap\_vecbench\_cross\_run.py}; full report at
\texttt{research/paper/vecbench\_cross\_run\_ci.md}). The point estimate next to
the CI is the p50 from the JSON cited in §5.2 (post-fix run, commit
\texttt{1e3d5f5}); the CI summarises run-to-run variance at this device +
harness preset. SQLiteAI rows do not carry a CI in this revision because the
precompute path's headline number (1 407 µs) comes from the dedicated
SQLite-family sweep, not the cross-engine run, and the dedicated 3-round audit
(3 072 ± 4 µs at the same N, companion §2 Table 2) is at a different operating
context --- see the \texttt{*} footnote below for the reconciliation. A future
revision will replace the cross-run bootstrap with a per-query bootstrap once
the new harness emits the full \texttt{latencies\_us} array (the field is added
to \texttt{VectorBenchmark.kt} in this revision; today's benches on Huawei P20
Lite, Huawei Y9 2019 and Moto G30 are already running with the patched harness
--- see §6.3 cross-platform status).

\(*\) The three SQLiteAI rows are \textbf{single-shot from the cross-engine run}
(the same paper384\_scale pass that produced every other row in this table). The
dedicated SQLite-family 3-round sweep at the same N = 20 000 reports
\texttt{default} 9 812 ± 363 µs, \texttt{optimized} 9 570 ± 548 µs,
\texttt{precompute} 3 072 ± 4 µs (companion report §2, Table 2; full audit at
\texttt{research/PAPER\_T11\_T12\_RECONCILIATION.md}). We keep the cross-engine
number here so every row in Table 11 is on the same measurement context
(single-shot, back-to-back schedule), and report the dedicated mean ± SD in the
companion. The qualitative comparison against Dazzle SQ8 (208 µs at the same
operating point) is the same direction in both readings --- Dazzle HNSW
\(O(\log N)\) stays inside the sub-millisecond envelope, the SIMD-accelerated
linear-scan path of SQLiteAI precompute does not. The exact numerical ratios at
this operating point, together with the Algorithm-class disclosure that frames
them, live in §5.8.4.

\(\dagger\) The Dazzle/Valkey rows run in-process and the bench did not exercise
an on-disk persistence path (BGSAVE was disabled to keep the I/O cost out of the
measurement); \texttt{INFO\ memory\ →\ used\_memory\_dataset} reads back as
\textasciitilde50 KB because the HNSW graph plus the quantised vector arena live
outside Valkey's accounted arenas (simsimd's aligned allocator). To put the
Dazzle rows on the same column as the SQLite-family \texttt{db\_file\_bytes}
numbers, we report the \textbf{analytical RAM/disk footprint} instead, computed
as \texttt{(N\ ×\ dim\ ×\ bytes\_per\_component)\ +\ (N\ ×\ M\ ×\ 4)} --- the
vectors plus the HNSW graph (int32 neighbour lists, \texttt{M\ =\ 32} per node).
At N = 20 000, dim = 384, M = 32: graph alone = 2.44 MB; SQ8 vectors = 7.32 MB →
\textbf{9.77 MB}; F16 = 14.65 MB → 17.09 MB; HNSW (fp32) = 29.30 MB → 31.74 MB;
SQ8+Rerank (int8 main + fp32 side-store) = 36.62 MB → 39.06 MB. These numbers
are conservative upper bounds in the sense that they ignore allocator alignment
slack but lower bounds in the sense that they do not include the per-node level
list overhead (which is small: \textasciitilde{}\(N \cdot \log_2 M\) bytes ≈
tens of KB at this scale).

The scientifically primary comparison is Dazzle vs ObjectBox because both are
HNSW implementations.

At recall floor \(\geq\) 0.95 (the common mobile RAG operating point), Dazzle
SQ8 reaches 208 µs p50 search at N = 20 000 vs ObjectBox 4.x at 853 µs (4.1×
faster), with both engines on HNSW. At strict recall \(\geq\) 0.99, ObjectBox
keeps a higher recall (0.994 vs 0.959 for Dazzle SQ8 at ef = 10) and Dazzle SQ8
sits at a lower-recall operating point in this configuration. This is an
expected trade-off between the int8 quantisation Dazzle SQ8 uses and the float32
precision ObjectBox keeps, not a contradiction.

SQLiteAI remains a useful secondary reference because it represents the current
SQLite production path in many mobile stacks. The observed ratio against
SQLiteAI at N = 20 000 is consistent with O(log N) vs O(N) behavior and should
be read with that algorithmic note attached.

\subsubsection{5.8.3 Envelope Interpretation}\label{envelope-interpretation}

Across the 9-configuration grid, Dazzle stays in the sub-millisecond search
envelope for the target mobile range and ObjectBox remains in the same order of
magnitude under matched HNSW settings. The relevant engineering conclusion is
that both systems are viable HNSW-class choices for on-device RAG, with
different constant-factor trade-offs by operating point and recall target.

Future versions of SQLiteAI sqlite-vector are expected to expose HNSW. At that
point the meaningful comparison will shift from algorithm class asymmetry to
constant-factor analysis between two HNSW implementations.

\subsubsection{5.8.4 SQLite Extension Variant Sweep ---
summary}\label{sqlite-extension-variant-sweep-summary}

\begin{quote}
\textbf{Algorithm-class disclosure.} This subsection compares Dazzle SQ8 (HNSW,
\(O(\log N)\) expected search) against SQLiteAI sqlite-vector v0.9.95
(SIMD-accelerated quantised linear scan, \(O(N)\)). The two systems are in
\textbf{different algorithm classes}, not the same class with different
constants. The numerical ratios reported below (e.g.~6.8×--14.8× at N = 20 000)
are therefore \emph{envelope references} --- they describe how far inside or
outside the sub-millisecond search envelope each path lands at a representative
mobile-RAG operating point --- and \textbf{not engine-quality claims} about
SQLiteAI as a product. The defensible same-class headline comparison in this
paper is \textbf{Dazzle SQ8 vs ObjectBox 4.x} in §5.8.2 (both HNSW). When a
future SQLiteAI release exposes HNSW, the meaningful comparison will shift to
constant-factor analysis between two HNSW implementations and the ratio framing
here will be obsolete; until then we report it because the algorithm-class
asymmetry is itself the engineering signal a mobile-RAG team needs at the
corpus-size threshold this section documents.
\end{quote}

To close the ``single SQLite path'' gap we ran a focused N-sweep on Moto G35 5G
over SQLite-family vector backends only --- \texttt{sqlite\_plain},
\texttt{sqlite\_vec} and SQLiteAI in their \texttt{default} / \texttt{optimized}
/ \texttt{precompute} variants --- at dim = 384, k = 10, 100 queries, 3 rounds,
N \(\in\) \{200, 1 k, 5 k, 10 k, 20 k\}. The N = 20 000 operating-point row for
each variant is merged into Table 11 above so the cross-engine comparison is in
one place; the per-N detail (retrieval p50, ingest total, on-disk footprint)
lives in \textbf{§2 of the companion engineering report}
(\texttt{research/paper/companion\_engineering\_report.md}, Tables 2 / 3 / 4).

The headline reading is \textbf{about envelope membership, not winner
declaration}. At N = 20 000, dim = 384, k = 10, recall ≥ 0.95:

\begin{itemize}
\tightlist
\item
  \textbf{Dazzle SQ8} lands at 208 µs (Table 11) --- clearly inside the
  sub-millisecond search envelope.
\item
  \texttt{sqlite\_vector\_ai\_precompute} lands at \textbf{1 407 µs} in the
  cross-engine single-shot (Table 11) and at \textbf{3 072 µs ± 4 µs} mean
  across the 3-round dedicated sweep (companion §2 Table 2; full audit at
  \texttt{research/PAPER\_T11\_T12\_RECONCILIATION.md}) --- outside the
  sub-millisecond envelope, into single-digit milliseconds.
\item
  \texttt{sqlite\_vec} (brute-force linear scan) is consistently above both
  SQLiteAI readings.
\end{itemize}

The 6.8×--14.8× factor between Dazzle SQ8 and SQLiteAI precompute at this N is
therefore \emph{evidence that the HNSW path stays inside the sub-millisecond
envelope while the linear-scan path crosses into the millisecond range} ---
i.e.~it documents a corpus-size threshold beyond which engine choice starts to
matter at the agent-loop level. For corpus sizes where every engine stays under
the inference noise floor (small N, e.g.~N = 200), the choice should default to
whichever engine has the primitive surface the application needs (Table 1, §2.2)
regardless of the latency delta. This comparison still spans two algorithm
classes (HNSW \citep{hnsw} \(O(\log N)\) vs linear scan \(O(N)\)); the exact
ratio is an envelope reference, not a claim of intrinsic engine superiority.

\subsection{End-to-End Application on Moto G35 5G: Small-Model + On-Device
RAG}\label{end-to-end-application-on-moto-g35-5g-small-model-on-device-rag}

\begin{quote}
\textbf{Scope of this section.} The end-to-end RAG ablation reported below runs
on a single physical device (Motorola Moto G35 5G, Snapdragon-class budget
Android handset). Cross-platform end-to-end validation across additional SoCs
(and an iPhone 12 Pro end-to-end run, currently storage-path only) is reported
in the next revision; storage-path validation on iPhone 12 Pro already appears
in §5.5, and all twelve backends in the §5.2 / §5.3 tables run on both
platforms. The §5.9 conclusions about retrieval-as-accuracy-lever therefore
stand on a single-device end-to-end run + cross-platform storage-path parity,
not on end-to-end coverage across hardware.
\end{quote}

This section's thesis is application-level: the availability of an in-process
vector primitive enables a qualitatively different agent pattern --- small-model
+ on-device RAG \citep{lewisrag, karpukhindpr} --- that can achieve higher
factual accuracy than a larger model without retrieval. The full pipeline runs
on-device, with no cloud and no external server, including on a \$150 Android
phone class. The 2 K passage corpus we index also keeps every retrieved passage
well within the small-model context window, so retrieval-precision artefacts
caused by long-context degradation \citep{liulostmiddle} are not on the path
between the index and the answer.

\subsubsection{5.9.1 Setup}\label{setup}

We index 2 000 passages from a deterministic mini-slice of Natural Questions
\citep{nq2019} on a Motorola Moto G35 5G (Unisoc T760, 4 GB RAM, Android 14).
Passages and (question, gold-passage) pairs come from the SBERT pair split of NQ
\citep{sbertnq}; canonical short-answer aliases are joined from
\texttt{nq\_open} \citep{nqopen, leeorqa} by lowercased question text.
\textbf{Alias precondition.} Not every NQ question has a canonical short answer
in \texttt{nq\_open} --- the join produces a \texttt{short\_answers} field only
for the subset of questions that the NQ-open authors annotated. We pre-filter
the candidate pool to keep only questions with at least one short-answer alias,
then sample 200 of those (so every query in the bench has a non-empty
ground-truth set for \texttt{EM\_short} / \texttt{F1\_short}). 2 000 passages
are constructed as the 200 gold positives plus 1 800 random distractors drawn
from the same pool, seed 42; passages are truncated to 1 800 characters to fit
\texttt{bge-small-en-v1.5}'s 512-token context. The slice is regenerable from
\texttt{research/scripts/nq\_slice.py} and is fingerprinted by sha256 prefix
\texttt{63be4b8894c71ff3} (full provenance in
\texttt{research/data/nq\_slice/README.md}). Embeddings produced with
\texttt{bge-small-en-v1.5} \citep{bgesmall} (q4\_k\_m, 24 MB, dim 384, n\_ctx
512; the BGE-small model carries an MTEB \citep{mteb} score of 62.2 on average
and ranks within the top-quartile of \textless100 MB models on the HF
leaderboard at submission time); Dazzle HNSW\_SQ8 index with M=32, efC=400. For
each of the 200 queries we retrieve the k=5 most cosine-similar passages with
efRuntime=64 and inject them into the prompt. Configurations:

\begin{itemize}
\tightlist
\item
  \textbf{Qwen 2.5 0.5B Instruct + Dazzle RAG.} GGUF q4\_k\_m, 380 MB on disk.
  n\_ctx=2 048. Prompt formed as
  \texttt{\textless{}system\textgreater{}…\ \textless{}retrieved\ passages\textgreater{}…\ \textless{}user\ question\textgreater{}}.
  Greedy decoding, max\_new\_tokens=64.
\item
  \textbf{Qwen 2.5 1.5B Instruct without RAG.} GGUF q4\_k\_m, 940 MB on disk ---
  2.5× larger, 3× more parameters. n\_ctx=2 048. Prompt formed as
  \texttt{\textless{}system\textgreater{}\ \textless{}user\ question\textgreater{}}.
  Greedy decoding, max\_new\_tokens=64.
\end{itemize}

\textbf{Evaluation protocol.} For each query the model generates up to 64
tokens. We extract the predicted answer span \texttt{y\_pred} as the prefix of
the generation up to the first newline or fresh \texttt{Question:} /
\texttt{Answer:} echo (the prompt already terminates in \texttt{Answer:}, so a
second occurrence is hallucinated continuation that we discard). Both
\texttt{y\_pred} and the gold short answer \texttt{y\_gold} are normalised
following the SQuAD v1.1 protocol \citep{squad} (lowercasing, removal of the
articles \texttt{\{a,\ an,\ the\}} and punctuation, whitespace collapsing). We
then report three short-answer metrics plus a passage backup:

\begin{itemize}
\tightlist
\item
  \textbf{\texttt{EM\_short}} (strict) --- \texttt{1} iff
  \texttt{norm(y\_pred)\ ==\ norm(y\_gold)} for any alias, \texttt{0} otherwise.
\item
  \textbf{\texttt{F1\_short}} --- token-level F1 (whitespace tokenisation)
  between \texttt{norm(y\_pred)} and the best-matching \texttt{norm(y\_gold)}
  alias, averaged over queries. By construction
  \texttt{EM\_short\ ≤\ F1\_short}.
\item
  \textbf{\texttt{EM\_contains}} (lax) --- \texttt{1} iff any alias's normalised
  tokens appear as a contiguous substring anywhere in the \emph{full} generation
  (not just in the extracted span). This separates \emph{``the model knows the
  answer''} from \emph{``the model answers concisely''}; it is reported as a
  diagnostic rather than the primary metric.
\item
  \textbf{\texttt{F1\_passage}} --- token-F1 between \texttt{y\_pred}
  (whitespace tokens, no aggressive normalisation) and the full gold passage; it
  measures whether the generation reproduces the supporting context, independent
  of conciseness.
\end{itemize}

Empty generations and refusals score \texttt{0} on every metric. Note that
\texttt{F1\_short} and \texttt{EM\_contains} are complementary signals, not
nested: a short span can have token overlap with a long alias yet fail the
contiguous substring test (e.g.~\texttt{y\_pred\ =\ "1957"},
\texttt{y\_gold\ =\ "1\ April\ 1957"} gives \texttt{F1\_short\ =\ 0.5} but
\texttt{EM\_contains\ =\ 0}).

\textbf{Scorer audit.} The scoring pipeline runs in two textually parallel
implementations: the on-device Kotlin scorer at
\texttt{experiment/backends/android/core/RagE2EBench.kt} (functions
\texttt{emStrict}, \texttt{emContains}, \texttt{f1Short},
\texttt{extractAnswerSpan}) and the offline Python re-scorer at
\texttt{research/scripts/recompute\_rag\_metrics.py} (same four entry points).
Both share the SQuAD v1.1 \citep{squad} normalisation contract (lowercase + drop
\texttt{\{a,\ an,\ the\}} + strip ASCII punctuation + collapse whitespace) and
the same span terminators (\texttt{\textbackslash{}n}, \texttt{Question:},
\texttt{Answer:}). The audit suite at
\texttt{research/scripts/test\_rag\_scoring.py} (31 unit tests, eight tiers:
normalisation, span extraction, \texttt{EM\_short}, \texttt{EM\_contains},
\texttt{F1\_short}, the two SQuAD invariants \texttt{EM\_short\ ≤\ F1\_short}
and \texttt{EM\_short\ ≤\ EM\_contains}, an 18-case cross-platform fixture at
\texttt{research/scripts/fixtures/rag\_scoring\_cases.json} that both
implementations consume, and a regression test that re-scores the JSON cited in
Table 15 and verifies the aggregates match within ±0.005 absolute) runs in the
\texttt{paper-consistency} GitHub Actions workflow on every push to
\texttt{main} and every PR that touches the paper, the scorer, the cited JSON,
or the fixture. The expanded \texttt{check\_paper\_consistency.py} carries two
dedicated checks for this audit (\texttt{SQuAD-scorer-audit},
\texttt{Kotlin-Python-scorer-parity}); 19/19 checks green at the time of this
revision.

\subsubsection{5.9.2 Result --- full 2×2 matrix}\label{result-full-22-matrix}

\textbf{Table 15 --- End-to-end RAG factual accuracy on 200 NQ queries} (Moto
G35 5G; greedy decoding, max\_new\_tokens = 64; same single bench run for all
four cells; raw JSON at
\texttt{research/benchmarks/results/Moto\_G35\_5G/rag\_2x2/rag\_e2e\_moto\_g35\_5G\_1777395311213.json},
SHA-256
\texttt{00d21f6c8752ffaa1015624b69a5e5d0fd403670d72561e3838bdac0ab461e76}).

Each cell is reported as point estimate with the bootstrap 95 \%
percentile-method confidence interval \texttt{{[}lo,\ hi{]}} over the per-query
metric array (n = 200, B = 10 000, seed = 42). The full bootstrap report
(per-cell CIs, paired-qid ratio CIs across all 16 metric × ratio cells,
significance flags) lives in \texttt{research/paper/rag\_2x2\_with\_ci.md} and
is regenerable from the same JSON via
\texttt{research/scripts/bootstrap\_rag\_metrics.py}.

{\def\LTcaptype{none} % do not increment counter
\begin{longtable}[]{@{}
  >{\raggedright\arraybackslash}p{(\linewidth - 10\tabcolsep) * \real{0.1929}}
  >{\raggedright\arraybackslash}p{(\linewidth - 10\tabcolsep) * \real{0.0643}}
  >{\raggedright\arraybackslash}p{(\linewidth - 10\tabcolsep) * \real{0.1857}}
  >{\raggedright\arraybackslash}p{(\linewidth - 10\tabcolsep) * \real{0.1857}}
  >{\raggedright\arraybackslash}p{(\linewidth - 10\tabcolsep) * \real{0.1857}}
  >{\raggedright\arraybackslash}p{(\linewidth - 10\tabcolsep) * \real{0.1857}}@{}}
\toprule\noalign{}
\begin{minipage}[b]{\linewidth}\raggedright
Configuration
\end{minipage} & \begin{minipage}[b]{\linewidth}\raggedright
Size
\end{minipage} & \begin{minipage}[b]{\linewidth}\raggedright
EM\_short
\end{minipage} & \begin{minipage}[b]{\linewidth}\raggedright
EM\_contains
\end{minipage} & \begin{minipage}[b]{\linewidth}\raggedright
F1\_short
\end{minipage} & \begin{minipage}[b]{\linewidth}\raggedright
F1\_passage
\end{minipage} \\
\midrule\noalign{}
\endhead
\bottomrule\noalign{}
\endlastfoot
Qwen 0.5B (no RAG) & 380 MB & 0.015 {[}0.000, 0.035{]} & 0.105 {[}0.065,
0.150{]} & 0.079 {[}0.055, 0.106{]} & 0.151 {[}0.138, 0.164{]} \\
Qwen 0.5B + Dazzle RAG & 380 MB & \textbf{0.120} {[}0.080, 0.165{]} &
\textbf{0.630} {[}0.565, 0.695{]} & \textbf{0.235} {[}0.191, 0.283{]} & 0.334
{[}0.300, 0.369{]} \\
Qwen 1.5B (no RAG) & 940 MB & 0.045 {[}0.020, 0.075{]} & 0.110 {[}0.070,
0.155{]} & 0.118 {[}0.084, 0.154{]} & 0.085 {[}0.073, 0.098{]} \\
Qwen 1.5B + Dazzle RAG & 940 MB & \textbf{0.310} {[}0.245, 0.375{]} &
\textbf{0.735} {[}0.670, 0.795{]} & \textbf{0.487} {[}0.430, 0.543{]} &
\textbf{0.235} {[}0.208, 0.265{]} \\
\end{longtable}
}

Reading the 2×2 across rows and columns (ratio CIs are bootstrap percentile
intervals on the \textbf{paired-qid} ratio statistic --- same draw of query
indices in numerator and denominator on each iteration; the full ratio table is
in \texttt{rag\_2x2\_with\_ci.md} Table 15b):

\begin{itemize}
\tightlist
\item
  \textbf{Retrieval helps both model sizes.} Adding Dazzle RAG lifts
  \texttt{EM\_contains} from 0.105 → 0.630 on the 0.5B (\textbf{6.00× {[}4.19×,
  9.71×{]}}) and 0.110 → 0.735 on the 1.5B (\textbf{6.68× {[}4.66×, 10.86×{]}});
  the strict \texttt{EM\_short} lift is 8.00× {[}3.00×, 29.00×{]} on the 0.5B
  (\texttt{⚠} --- denominator base rate 0.015 makes the multiplicative ratio
  statistically unstable; the additive lift 0.120 − 0.015 = 0.105 {[}bootstrap
  95 \% CI approximately ±0.04{]} is the better summary at this base rate) and
  6.89× {[}4.00×, 16.00×{]} on the 1.5B. Every reported ratio's 95 \% CI
  excludes 1.0 --- every directional effect in the table is significant at the
  conventional bootstrap-percentile level. The retrieval gain is \textbf{larger
  in absolute terms on the bigger model} --- a 1.5B model with a retrieval
  primitive in reach uses the retrieved context more effectively than a 0.5B
  model does. This contradicts a naive reading where retrieval only matters for
  small models.
\item
  \textbf{Small + RAG outperforms large no-RAG on every factual metric.} Qwen
  0.5B + RAG reaches \texttt{EM\_short} 0.120 / \texttt{EM\_contains} 0.630 vs
  Qwen 1.5B (no RAG) 0.045 / 0.110 (the cross-cell ratio on
  \texttt{EM\_contains} is \textbf{5.73× {[}4.07×, 9.00×{]}}, paired-qid) ---
  i.e.~a 380 MB model with retrieval beats a 3× larger 940 MB model without it,
  confirming the v1 reading at the lower-bound corner of the matrix even with
  the corrected SQuAD-normalised metrics.
\item
  \textbf{Large + RAG is the global maximum on every factual metric.} Qwen 1.5B
  + RAG reaches \texttt{EM\_short} 0.310 / \texttt{EM\_contains} 0.735 ---
  \textbf{2.58× {[}1.82×, 4.06×{]}} and \textbf{1.17× {[}1.05×, 1.30×{]}} over
  Qwen 0.5B + RAG respectively. The accuracy gain comes at a 2.8× per-turn
  latency cost (49.23 s vs 17.62 s p50, Table 16), which makes this operating
  point appropriate for batch or asynchronous workloads rather than interactive
  turns (see §5.9.4); the storage backend's contribution to that turn budget is
  a constant \textasciitilde22 ms + 0.6 ms regardless of which decoder consumes
  the result.
\end{itemize}

The three short-answer metrics tell consistent but distinct stories.
\texttt{EM\_contains} shows that retrieval surfaces the right \emph{facts}, even
when the model wraps them in extra prose. \texttt{F1\_short} measures how close
the \emph{first-line span} comes to the gold form when normalised --- partial
matches like \texttt{y\_pred\ =\ "Mark\ Skaife\ won\ the\ championship"} against
\texttt{y\_gold\ =\ "Mark\ Skaife"} score around 0.5. \texttt{EM\_short} is the
strictest reading: the extracted span, after normalisation, must equal the gold
answer exactly. The invariant \texttt{EM\_short\ ≤\ F1\_short} and
\texttt{EM\_short\ ≤\ EM\_contains} holds in all four rows, as required by the
definitions in §5.9.1.

\textbf{Erratum --- internal v1 of this section.} An earlier internal-only draft
of §5.9.2 (commit \texttt{f999003}, never published to arXiv or any public
preprint server) reported buggy short-answer metrics. \texttt{EM\_short} was
implemented as a substring match on the full 64-token generation (so a verbose
but correct reply scored 1), and \texttt{F1\_short} was computed token-by-token
over that same full generation (so the same verbose reply was heavily penalised
by all the surrounding prose). Together they produced the pair
\texttt{EM\_short\ =\ 0.665}, \texttt{F1\_short\ =\ 0.077} for small+RAG, which
is \textbf{mathematically impossible} under any SQuAD-style definition
(\texttt{EM\ ≤\ F1} per sample, hence per average; see \citep{squad} for the
canonical definitions). Both formulas were wrong, not a different methodological
choice. The numbers in the table above come from a fresh 200-query run with the
corrected scoring as described in §5.9.1.

\textbf{The v1 and post-fix runs are two related but not byte-identical
experiments.} A direct generation-level diff between the two runs on the 10 (qid
× variant) pairs that v1 persisted to its JSON (the v1 harness wrote only the
first 10 of the 200-query run to \texttt{examples}; the post-fix harness writes
all 200) shows that the model artefacts changed between the two runs:

\begin{itemize}
\tightlist
\item
  \texttt{small\_llm} \texttt{qwen2.5-0.5b-instruct-q4\_k\_m.gguf} grew from
  397,808,192 B (v1) to 491,400,032 B (post-fix), a +94 MB delta;
\item
  \texttt{large\_llm} \texttt{qwen2.5-1.5b-instruct-q4\_k\_m.gguf} grew from
  986,048,768 B (v1) to 1,117,320,736 B (post-fix), a +131 MB delta;
\item
  the embedder \texttt{bge-small-en-v1.5-q4\_k\_m.gguf} is byte-identical in
  both runs (24,808,576 B).
\end{itemize}

Both LLM files share the same upstream Qwen 2.5 base weights but are different
\texttt{q4\_k\_m} repacks, so the on-device generations differ even on identical
prompts: 0 of 10 anchor generations are identical in \texttt{small\_rag}, 1 of
10 in \texttt{large\_no\_rag}. The v1 → post-fix \texttt{EM\_short} /
\texttt{F1\_short} delta is therefore a \textbf{combined effect} of (a) the
scorer-formula fix described in the previous paragraph and (b) a model-artefact
change between runs, not a pure scorer-isolation diff. The full per-row
verifiable diff is at \texttt{research/results/v1\_vs\_postfix\_diff.md},
generated by \texttt{research/scripts/diff\_v1\_postfix\_rag.py} from the v1
JSON recovered out of git history (commit \texttt{80f465e}) and the post-fix
JSON cited in Table 15.

The \texttt{F1\_passage} column nevertheless lands close in both runs: v1
internal draft 0.385 (small+RAG) / 0.074 (large no-RAG) at n = 200 vs post-fix
0.334 / 0.085 (Table 15) --- i.e.~a ±0.05 shift that is consistent with two
related q4\_k\_m repacks of the same base model on the same workload, not an
order-of-magnitude change. We flag this here so any reviewer who saw the v1
draft can reconcile the numbers honestly: the canonical run for Table 15 is the
post-fix one cited in the table. The pre-fix output JSONs
(\texttt{research/results/rag\_e2e\_moto\_g35\_5G.json},
\texttt{research/results/rag\_e2e\_moto\_g35\_5G\_em.json},
\texttt{research/results/rag\_e2e\_moto\_g35\_5G\_em\_q200.json}) were retracted
from the artifact in commit \texttt{2802527} (see
\texttt{research/results/RAG\_RETRACTION.md}); the v1 source-code commit
\texttt{f999003} remains in the public git history for auditability, but its
\texttt{EM\_short} and \texttt{F1\_short} values must be discarded. The post-fix
replacement is the single 2×2 bench run cited in Table 15.

\subsubsection{5.9.3 Latency and Trade-Off
Space}\label{latency-and-trade-off-space}

\textbf{Table 16 --- End-to-end RAG turn latency breakdown} (Moto G35 5G, p50
over 200 queries; same 2×2 matrix as Table 15).

{\def\LTcaptype{none} % do not increment counter
\begin{longtable}[]{@{}lrrrrr@{}}
\toprule\noalign{}
Configuration & Embed & Search & Prefill & Decode & Total p50 \\
\midrule\noalign{}
\endhead
\bottomrule\noalign{}
\endlastfoot
Qwen 0.5B (no RAG) & --- & --- & 0.68 s & 0.97 s & 1.73 s \\
Qwen 0.5B + Dazzle RAG & 21.5 ms & 614 µs & 15.03 s & 2.92 s & 17.62 s \\
Qwen 1.5B (no RAG) & --- & --- & 2.01 s & 0.97 s & 2.98 s \\
Qwen 1.5B + Dazzle RAG & 23.9 ms & 692 µs & 45.89 s & 5.70 s & 49.23 s \\
\end{longtable}
}

The retrieval-enabled turn is \textbf{10× slower} for the 0.5B model (17.62 s vs
1.73 s) and \textbf{17× slower} for the 1.5B model (49.23 s vs 2.98 s). The
dominant cost is prefill, not vector lookup: embedding is 21--24 ms and search
is 614--692 µs across both RAG rows, while prompt length grows from
\textasciitilde27 tokens (no RAG) to \textasciitilde570 tokens (with retrieved
passages), and prefill scales with prompt length. Vector search across the four
cells stays within ±15 \% of the 622 µs the §5.8 envelope predicted, so the
search component is \textbf{not} sensitive to which decoder is consuming the
result.

This 2×2 should therefore be read as a trade-off map, not a victory statement.
The latency cost of RAG is paid almost entirely in prefill, which is a property
of the LLM-decoder choice (1.5B or 0.5B), not of the storage backend.
Applications that prioritise factual grounding accept the extra prefill time;
applications that prioritise short-turn interaction can step down to a smaller
model and rely on retrieval to recover the factual grounding (e.g. 0.5B + RAG at
17.6 s p50 still beats 1.5B no-RAG at 2.98 s p50 on every accuracy metric in
Table 15). The key point is that the \textbf{RAG pattern is available on-device
with no cloud dependency}; the storage backend's contribution to that turn
budget is the \textasciitilde22 ms + 0.6 ms embed-and-search cost --- under 0.05
\% of the fastest RAG row's total turn.

\subsubsection{5.9.4 Implications for On-Device
Deployment}\label{implications-for-on-device-deployment}

Four practical consequences follow from the 2×2:

\begin{enumerate}
\def\labelenumi{\arabic{enumi}.}
\tightlist
\item
  \textbf{Retrieval, not size, is the dominant accuracy lever at this scale.}
  Both no-RAG cells sit at \texttt{EM\_contains} ≤ 0.110; adding Dazzle RAG
  lifts the 0.5B row to 0.630 (6.0×) and the 1.5B row to 0.735 (6.7×). The
  retrieval column-effect (≥6×) is almost an order of magnitude larger than the
  model-size row-effect at matched retrieval (1.5B + RAG / 0.5B + RAG = 1.17× on
  \texttt{EM\_contains}). For a factual agent on a phone, the first question is
  ``is the retrieval primitive available,'' not ``which weights fit in RAM.''
\item
  \textbf{A 380 MB model with retrieval beats a 940 MB model without retrieval
  on every factual metric in Table 15.} Stepping down from 1.5B to 0.5B
  parameters --- saving 560 MB on disk and proportional RAM --- is a net
  accuracy win when the agent has access to a retrieval layer, and an 8.7×
  latency improvement on the RAG row (49.23 s → 17.62 s p50, Table 16). For a
  memory-constrained Android device the practical recommendation is 0.5B +
  Dazzle RAG.
\item
  \textbf{The 1.5B + RAG cell is the global maximum}, but its 49.23 s turn
  latency makes it appropriate for batch / asynchronous workloads (overnight
  indexing, background summarisation) rather than interactive turns. The
  frontier is not a single best configuration --- it is a Pareto curve where
  0.5B + RAG and 1.5B + RAG occupy the latency-vs-accuracy extremes.
\item
  \textbf{Retrieval cost is negligible.} Vector search across both RAG rows is
  614--692 µs over 2 000 passages --- between 0.0014 \% (1.5B + RAG) and 0.0035
  \% (0.5B + RAG) of total turn latency. The embedder (21.5--23.9 ms p50) is the
  only storage-side component that scales with corpus size; in production with a
  static corpus it runs only at ingest, and a smaller embedder
  (\texttt{bge-micro}, dim 384, \textasciitilde5 MB) could cut that latency
  4--5× with no appreciable recall loss --- future work.
\end{enumerate}

This experiment compares ``with retrieval'' vs ``without retrieval.'' It is not
a direct application-level horse race between Dazzle, ObjectBox, and SQLiteAI
vector engines. We did not run this exact end-to-end experiment with alternative
vector backends due to engineering cost; we expect \texttt{EM\_contains} to
remain within approximately ±2 \% across backends at matched recall settings,
because retrieval recall dominates this metric while retrieval latency (614--692
µs) is \textless{} 0.005 \% of total turn latency on both RAG rows.

\section{Discussion}\label{discussion}

\subsection{Materialized Aggregates as a Primitive, Not an
Optimization}\label{materialized-aggregates-as-a-primitive-not-an-optimization}

The key mechanism behind low-latency retrieval in this paper is the
materialized-aggregates pattern, not a specific brand of database engine.
Whether implemented as a Lua-maintained hash in Dazzle, a trigger-maintained
table in SQLite, or an explicit update step in an in-memory structure, the
pattern shifts aggregation work to write time and makes read-time retrieval
bounded.

Dazzle's contribution is not inventing this pattern. The contribution is
shipping it as a first-class primitive surface (\texttt{EVALSHA}-maintained hash
+ \texttt{HMGET}) with typed SDK access, so the application author does not need
to build migration logic, trigger orchestration, or manual maintenance code for
each app.

This distinction matters to the paper's thesis: in the latency envelope where
all backends are acceptable, the deciding factor is less ``can the pattern
exist?'' and more ``how much engineering ceremony is required to make it correct
and maintainable in production.''

\subsection{The Bounded Context Block}\label{the-bounded-context-block}

The \textasciitilde58--62 token property independent of N is a consequence of an
architectural choice, not a coincidence. The context block encodes six
aggregated scalars (min/max/avg/count/\ldots) plus the probabilistic summary ---
not a list of N raw readings. This decouples prompt size from dataset size. An
agent running for weeks on a device should not progressively consume more of the
model's context window simply because more data has accumulated.

\subsection{Limitations}\label{limitations}

\begin{itemize}
\item
  \textbf{Primitive surface counted as native API, not as expressive
  equivalence.} Table 1 (§2.2) reports primitive availability as a categorical
  check (\texttt{N} native API / \texttt{E} official extension / \texttt{A}
  app-level reimplementation / \texttt{T} third-party extension) on how the
  backend reaches each construct, with the headline count measured on \texttt{N}
  only. Primitives marked \texttt{E}, \texttt{A} or \texttt{T} for non-Dazzle
  backends are \emph{implementable} atomically through the underlying engine's
  general-purpose facilities --- for example, sorted-set range queries via
  \texttt{ORDER\ BY} + \texttt{WHERE\ BETWEEN} in SQLite, atomic float
  increments via \texttt{UPDATE\ …\ SET\ v\ =\ v\ +\ ?} inside a transaction,
  geographic indexing via SQLite's \texttt{R*Tree} extension (an official
  extension shipped with the SQLite amalgamation since 2008), bounded
  \texttt{MAXLEN} streams via triggers, and TTL via a scheduled \texttt{DELETE}
  job. The signal of Table 1 is therefore ``how much glue code does the agent
  author write to reach this primitive'' (a primitive-surface ergonomic question
  that ties back to Table 9), not ``is this functionality achievable in this
  engine'' --- which is, in almost every case, yes. A reviewer reading Table 1
  as the latter is reading the wrong axis; the LOC differential in Table 9 is
  the more direct measurement of the same phenomenon.
\item
  \textbf{Harness post-mortem.} A bug in the \texttt{SqliteBruteforceVector}
  truth-source caused stale data to accumulate across configurations of the
  recall-floor benchmark in v1; it was fixed in commit \texttt{1e3d5f5} and
  every cell in Tables 4, 11 and 12 of this revision comes from the post-fix
  run. The full timeline (introduction, detection, fix, asserts that catch a
  future regression, and an audit of every other on-device benchmark in the repo
  for the same bug pattern) is documented in Appendix B
  (\texttt{research/paper/appendix\_b\_harness\_postmortem.md}). The audit found
  no other latent instance of the pattern at the time of this revision.
\item
  End-to-end LLM evaluation on a single physical device (Moto G35 5G) +
  cross-platform validation on iPhone 12 Pro. Results on other SoCs and OS
  versions may differ; we report results on three additional Android SoC
  families (Unisoc, MediaTek, Qualcomm) in companion benchmark reports.
\item
  Synthetic workload. Real IoT streams have bursty patterns, sensor failures,
  and additional concurrency we do not model.
\item
  Model accuracy evaluation uses three verifiable statistics; a more rigorous
  evaluation would use a larger factual question set and multiple runs per
  condition.
\item
  Incomplete iOS/Android transport parity: the SPSC ring buffer and the
  \texttt{io\_uring} path are Linux-only; iOS stays on the Phase-0 pipe path,
  which is sufficient for the evaluated workload (§5.5). First-class iOS parity
  is future work.
\item
  \textbf{Single-device end-to-end LLM workload.} The end-to-end RAG bench of
  §5.9 is reported on a single physical device (Moto G35 5G). The storage-only
  retrieval path is cross-validated on a physical iPhone 12 Pro and shows parity
  with the Moto G35 5G at the snapshot-cache layer (§5.5); cross-platform E2E
  LLM validation remains future work.
\item
  \textbf{Commercial peers evolve.} SQLiteAI v0.9.95 does not expose HNSW at the
  close of this bench; a future product version that includes HNSW would shift
  the §5.8 comparison from ``HNSW vs brute-force scan'' to ``HNSW vs HNSW'',
  reducing the gap to constant factors rather than algorithmic asymmetry. The
  bench pins a specific version (\texttt{vector-apple-xcframework-0.9.95.zip},
  SHA-256 verified) for reproducibility.
\item
  \textbf{Platform-allocator footprint skew.} The on-disk and in-memory size of
  an embedded backend is sensitive to the host allocator and to default
  configuration. We observed a 23× cross-platform delta on Dazzle (153 KB on iOS
  vs 6.7 KB on Android, same in- process Valkey-fork, same 200-reading load)
  attributable to the Apple \texttt{libsystem\ malloc} versus Android
  \texttt{jemalloc} arena behavior, and a 42× delta on RocksDB in the opposite
  direction (4.41 MB Android vs 104 KB iOS) driven by Android's WAL + SST
  pre-allocation defaults. Even with these skews the §5.3 comparison against
  SQLite still favours Dazzle by 24-56× on the same device, but the absolute
  single-platform numbers should be read as best-effort defaults, not tuned
  floors. Allocator-tuned footprint floors are future work.
\item
  \textbf{Default-configuration comparison in §5.2.} The storage-latency
  comparison uses default backend configurations. A SQLite setup with manually
  maintained materialized views via triggers would close most of the measured
  retrieval gap. Dazzle's contribution is shipping the pattern natively, not
  claiming algorithmic superiority over SQLite.
\item
  \textbf{Vector benchmark spans two algorithm classes.} In §5.8, Dazzle and
  ObjectBox run HNSW (O(log N)), while SQLiteAI v0.9.95 exposes brute-force
  quantized scan (O(N)). A future SQLiteAI release with user-accessible HNSW
  would shift the comparison to constant factors.
\item
  \textbf{RAG application benchmark isolates retrieval availability, not vector
  engine identity.} Section §5.9 compares with-retrieval vs without-retrieval
  and does not isolate Dazzle's vector engine against ObjectBox or SQLiteAI in
  the same end-to-end pipeline.
\item
  \textbf{Author bias on the LOC measurement.} The author of this paper is also
  the author of all six backend implementations against which Dazzle is measured
  in Table 9 (\texttt{SqliteContextManager}, \texttt{LmdbContextManager},
  \texttt{RocksDbContextManager}, \texttt{ObjectBoxContextManager},
  \texttt{InMemoryContextManager}, plus the Dazzle managers). The Dazzle-side
  managers were the first written and the most iterated; the
  SQLite/LMDB/RocksDB/ObjectBox managers are direct ports of the same agent
  state-layer surface but were not subsequently optimised for line count. A
  reviewer who treats the LOC numbers as a \emph{primitive surface} proxy (how
  much of the agent state-layer is expressed in domain-typed commands vs in glue
  code over a lower-level API) is reading them correctly; a reviewer who treats
  them as a \emph{code-quality} proxy is not.
\item
  \textbf{Conflict of interest and public commitment.} The author of this paper
  is also the author of \texttt{dazzle-sdk} (Apache-2.0; Maven Central / pub.dev
  / npm / Swift Package Manager) and of every comparison-engine wrapper
  benchmarked here
  (\texttt{experiment/backends/\{android,ios\}/\{sqlite,\ lmdb,\ rocksdb,\ objectbox,\ inmemory,\ dazzle\}/}).
  To make the comparison falsifiable rather than self-serving, three concrete
  commitments hold from the date of this revision forward:

  \begin{enumerate}
  \def\labelenumi{\arabic{enumi}.}
  \tightlist
  \item
    \textbf{All six backend wrappers are first-class artefacts in the
    repository.} Their full source, build configuration, and test entry points
    live under \texttt{experiment/backends/} and have the same Apache-2.0
    licence as Dazzle. Anyone can fork, re-tune, and re-run.
  \item
    \textbf{External pull requests that reduce the LOC count of any non-Dazzle
    wrapper, or that improve any non-Dazzle backend's measured numbers in Tables
    3 / 5 / 11, will be reviewed and merged on technical merit.} The LOC-9
    differential (175/186 for Dazzle vs 200--290 for the alternatives) is
    therefore an explicitly contestable claim, not a fixed verdict.
  \item
    \textbf{Every accepted optimisation is recorded in
    \texttt{experiment/backends/CHANGELOG.md} before the next benchmark freeze.}
    The CHANGELOG is part of the repository and acts as a public ledger of
    ``what improved which engine, and which paper revision incorporates it.'' A
    reviewer can audit the CHANGELOG between revisions to verify that
    contributed optimisations propagate into the paper's measurements.
  \end{enumerate}

  The benchmark methodology, raw JSON outputs, and analysis scripts are open in
  the repository (paths listed in §8.1 Reproducibility) so any claim in the
  paper can be re-run against an independent build of the comparison engines.
\end{itemize}

\subsubsection{Cross-platform validation of the vector
benchmark}\label{cross-platform-validation-of-the-vector-benchmark}

Tables 11 and 12 are reported on a single physical device (Moto G35 5G, Unisoc
T760, ARMv8.2-A). To rule out single-device artefacts, we re-ran the same
harness (same query set, same N ∈ \{200, 1k, 5k, 20k\}, same \texttt{k\ =\ 10},
dim = 384) on three additional Android SoC families and re-bootstrapped the
headline N = 20 000 cell with paired-query resampling (\texttt{B\ =\ 10\ 000},
seed = 42). The full per-engine tables are in
\texttt{research/paper/vecbench\_cross\_platform\_ci.md}; the headline ratios
are reproduced here:

{\def\LTcaptype{none} % do not increment counter
\begin{longtable}[]{@{}
  >{\raggedright\arraybackslash}p{(\linewidth - 8\tabcolsep) * \real{0.2579}}
  >{\raggedright\arraybackslash}p{(\linewidth - 8\tabcolsep) * \real{0.2390}}
  >{\raggedright\arraybackslash}p{(\linewidth - 8\tabcolsep) * \real{0.1384}}
  >{\raggedright\arraybackslash}p{(\linewidth - 8\tabcolsep) * \real{0.1824}}
  >{\raggedright\arraybackslash}p{(\linewidth - 8\tabcolsep) * \real{0.1824}}@{}}
\toprule\noalign{}
\begin{minipage}[b]{\linewidth}\raggedright
Device
\end{minipage} & \begin{minipage}[b]{\linewidth}\raggedright
SoC (ISA)
\end{minipage} & \begin{minipage}[b]{\linewidth}\raggedright
Dazzle SQ8 p50 µs
\end{minipage} & \begin{minipage}[b]{\linewidth}\raggedright
vs ObjectBox p50
\end{minipage} & \begin{minipage}[b]{\linewidth}\raggedright
vs SQLiteAI precompute p50
\end{minipage} \\
\midrule\noalign{}
\endhead
\bottomrule\noalign{}
\endlastfoot
Moto G35 5G \emph{(reference, Table 11)} & Unisoc T760 (ARMv8.2-A + fp16/dot) &
306 {[}296, 314{]} & 0.34× {[}0.32×, 0.36×{]} ★ & 0.10× {[}0.10×, 0.10×{]} ★ \\
Moto G30 & Snapdragon 662 (ARMv8.0-A) & 568 {[}530, 608{]} & 0.33× {[}0.31×,
0.35×{]} ★ & 0.11× {[}0.10×, 0.11×{]} ★ \\
Huawei Y9a (FRL-L23) & MediaTek Helio G80 (ARMv8.2-A) & 744 {[}718, 778{]} &
0.22× {[}0.21×, 0.23×{]} ★ & 0.15× {[}0.15×, 0.16×{]} ★ \\
Huawei P20 Lite (ANE-LX3) & Kirin 659 (ARMv8.0-A, Linux 4.9) & 974 {[}949,
991{]} & 0.30× {[}0.28×, 0.30×{]} ★ & 0.13× {[}0.13×, 0.14×{]} ★ \\
\end{longtable}
}

All four ratios are statistically significant (95 \% CI excludes 1.0); a Dazzle
SQ8 search is 3.0×--4.5× faster than ObjectBox 4.x HNSW and 6.7×--10.0× faster
than SQLiteAI v0.9.95 precompute on every chip we ran. Two findings to flag:

\begin{enumerate}
\def\labelenumi{\arabic{enumi}.}
\item
  \textbf{Cross-platform correctness required a single-binary build path on
  Helio G80 (FRL-L23) and Kirin 659 (ANE-LX3).} The reference Moto G35 binary
  used \texttt{-march=armv8.2-a+fp16+dotprod} and SIGILL-ed on ARMv8.0 chips at
  the first \texttt{FT.CREATE}. We retargeted Dazzle to
  \texttt{-march=armv8-a\ -mcpu=generic} with runtime SIMD dispatch via
  \texttt{simsimd}, disabled PAC/BTI in user code (the Kirin 659 kernel on EMUI
  8.2 is ≤ 4.9 and does not trap-and-emulate
  \texttt{MRS\ Xt,\ ID\_AA64ISAR0\_EL1} reads from EL0), and shipped a userspace
  \texttt{SIGILL}/MRS-emulation handler that reports ``no extensions'' to
  libraries that probe by direct MRS. This path is what landed in
  \texttt{f7a9832\ →\ cc6e731}.

  We then implemented a \emph{multi-target build} that ships \textbf{both} code
  paths in the same APK (\texttt{libdazzle.so},
  \texttt{-march=armv8-a\ -mcpu=generic}, baseline, 8.6 MB;
  \texttt{libdazzle\_v82.so},
  \texttt{-march=armv8.2-a+fp16+dotprod\ -mcpu=cortex-a78}, 8.7 MB) and
  dispatches at \texttt{System.loadLibrary} time on \texttt{/proc/cpuinfo}
  \texttt{asimdhp\ +\ asimddp} features (via
  \texttt{dev.dazzle.sdk.DazzleNativeLoader}). To rule out review criticism
  around ``comparison across SoCs is not apples-to- apples when each chip runs a
  different binary'', we ran the benchmark twice on every device (same harness,
  same 100-query set, same paired bootstrap):

  \begin{itemize}
  \tightlist
  \item
    \textbf{Round 1 --- dispatched/optimal}: each chip loads the variant its
    silicon can run at hardware speed (\texttt{libdazzle\_v82.so} on Unisoc T760
    + Helio G80, \texttt{libdazzle.so} baseline on Snapdragon 662 + Kirin 659).
  \item
    \textbf{Round 2 --- baseline-forced apples-to-apples}: every chip forces
    \texttt{libdazzle.so} baseline via
    \texttt{-\/-es\ force\_native\_variant\ baseline}, so the comparison across
    SoCs is exactly the same binary on every device.
  \end{itemize}

  The dispatched p50 (round 1) and the baseline-forced p50 (round 2) for
  \texttt{dazzle\_sq8} at N = 20 000 are:

  {\def\LTcaptype{none} % do not increment counter
  \begin{longtable}[]{@{}
    >{\raggedright\arraybackslash}p{(\linewidth - 8\tabcolsep) * \real{0.2013}}
    >{\raggedright\arraybackslash}p{(\linewidth - 8\tabcolsep) * \real{0.3377}}
    >{\raggedleft\arraybackslash}p{(\linewidth - 8\tabcolsep) * \real{0.1688}}
    >{\raggedleft\arraybackslash}p{(\linewidth - 8\tabcolsep) * \real{0.2013}}
    >{\raggedright\arraybackslash}p{(\linewidth - 8\tabcolsep) * \real{0.0909}}@{}}
  \toprule\noalign{}
  \begin{minipage}[b]{\linewidth}\raggedright
  Chip
  \end{minipage} & \begin{minipage}[b]{\linewidth}\raggedright
  SoC (ISA, big core)
  \end{minipage} & \begin{minipage}[b]{\linewidth}\raggedleft
  Round 1 p50 (dispatched)
  \end{minipage} & \begin{minipage}[b]{\linewidth}\raggedleft
  Round 2 p50 (baseline-forced)
  \end{minipage} & \begin{minipage}[b]{\linewidth}\raggedright
  Δ
  \end{minipage} \\
  \midrule\noalign{}
  \endhead
  \bottomrule\noalign{}
  \endlastfoot
  Moto G35 5G \emph{(reference)} & Unisoc T760 (ARMv8.2-A + fp16/dot,
  Cortex-A76) & 269 {[}261, 273{]} µs & 269 {[}264, 273{]} µs & \textbf{0 µs} \\
  Huawei Y9a (FRL-L23) & MediaTek Helio G80 (ARMv8.2-A + fp16/dot, A75) & 671
  {[}642, 701{]} µs & 588 {[}566, 612{]} µs & \textbf{-83 µs} ² \\
  Moto G30 & Snapdragon 662 (ARMv8.0-A, Cortex-A73) & 445 {[}409, 478{]} µs &
  519 {[}488, 548{]} µs & +74 µs ¹ \\
  Huawei P20 Lite (ANE-LX3) & Kirin 659 (ARMv8.0-A, Linux 4.9, Cortex-A53) &
  1054 {[}1035, 1092{]} µs & 1043 {[}1014, 1078{]} µs & -11 µs ¹ \\
  \end{longtable}
  }

  ¹ ARMv8.0 chips run the same \texttt{libdazzle.so} baseline in both rounds
  (the v82 variant would SIGILL); the round-to-round Δ is therefore pure
  thermal-/load-induced noise, ranging ±10 \% at N = 20 000 --- useful as a
  calibration of the per-run variance the dispatched-vs-baseline cell must clear
  to count as a real speedup.

  ² On Helio G80 (Cortex-A75 \citep{armcortexa75} big core) the dispatched v82
  binary is \textbf{slower} than the baseline by ≈ 12 \% at p50 (95 \% CIs do
  not overlap: {[}642, 701{]} vs {[}566, 612{]} µs). The hot-path NEON kernels
  are identical --- \texttt{simsimd} runtime- selects the same
  \texttt{\_neon\_f16\ /\ \_neon\_dotprod} symbols on both binaries --- so the
  gap is in the surrounding C++ code that the C++17 compiler emits for the
  wrapper around \texttt{simsimd\_*}.

  We initially attributed the regression to \texttt{-mcpu=cortex-a78}
  \citep{armcortexa78} issuing OoO speculation the A75 retire-port couldn't
  sustain, and ran a follow-up experiment (round 3): we recompiled
  \texttt{libdazzle\_v82.so} with \texttt{-mcpu=cortex-a76} \citep{armcortexa76}
  (one pipeline generation ahead of A75 instead of three) and re-benched FRL-L23
  in dispatched mode. Result: dazzle\_sq8 N = 20 000 p50 = \textbf{680 µs},
  indistinguishable from the round-1 cortex-a78 number (671 µs) and still well
  above the 588 µs baseline. The raw JSON and a one-page write-up of the
  experiment live at
  \texttt{research/benchmarks/results/multitarget/round3\_cortex\_a76\_test/}.

  The \texttt{-mcpu} retune therefore did \textbf{not} explain the regression;
  the scheduler model is not on the critical path. The remaining suspects are
  (a) \texttt{-march=armv8.2-a\ +\ fp16\ +\ dotprod} itself relaxing constraints
  the A75 microarchitecture cannot exploit (the C++17 compiler emits FP16
  truncate / extend chains and dotprod-style integer accumulators in the
  surrounding C++ wrapper code that chain-stall on A75's narrower issue ports),
  or (b) some inlined std::priority\_queue path that the wider target instructs
  the compiler to schedule more aggressively. Disambiguating those would require
  objdump-level diffing of the v82 vs baseline binaries on the few translation
  units in \texttt{valkeysearch\_module.cc}, which is out of scope for this
  revision.

  The shipped mitigation is the \texttt{force\_native\_variant} override the
  multi-target build already exposes: end users on Helio G80 / Cortex-A75 chips
  can opt into the baseline binary at runtime via
  \texttt{System.setProperty("dazzle.force\_native\_variant",\ "baseline")}
  before any Dazzle SDK call. A future revision could automate that by
  augmenting the \texttt{asimdhp\ +\ asimddp} heuristic with an
  \texttt{ro.board.platform} lookup that selects baseline when the chip is
  known-A75 (Helio G80 / G88, P60 etc.). The loader exposes both knobs already;
  we leave the lookup table empty in this revision because we have exactly one
  A75-class data point and a whitelist on n=1 would over-fit.

  \textbf{Headline finding (and a course-correction the reviewer should be aware
  of)}: on the only chip × engine cell where we can directly compare v82 against
  baseline on the same silicon (Moto G35 5G), the two binaries are
  \textbf{statistically indistinguishable} --- 269 µs vs 269 µs at N = 20 000,
  with overlapping 95 \% CIs. The ``31 \% perdido'' we initially attributed to
  the cross-platform \texttt{armv8-a} retarget (208 µs single-binary armv8.2
  build vs 306 µs cross-platform armv8-a build, both pre-multi-target) does not
  reproduce when we re-bench with the multi-target build: the dispatched binary
  recovers ≈ 12 \% vs the older armv8-a build (306 → 269 µs), not 31 \%. The
  remaining gap was either (a) measurement noise across runs (the ±10 \%
  round-to-round variance above is consistent with this), or

  \begin{enumerate}
  \def\labelenumii{(\alph{enumii})}
  \setcounter{enumii}{1}
  \tightlist
  \item
    compiler / scheduling differences in code outside the vector-search hot path
    (e.g.~ggml/llama.cpp baseline scaling the rest of the harness). Either way,
    the \emph{vector-search} hot path on a v8.2 chip is \textbf{already
    saturated} by \texttt{simsimd}'s runtime dispatcher: when the chip
    advertises \texttt{asimdhp\ +\ asimddp} in \texttt{/proc/cpuinfo}, simsimd
    selects the \texttt{*\_neon\_f16} / \texttt{*\_neon\_dotprod} kernels at
    first call and the per-kernel codegen is identical between the two binaries.
    The \texttt{-march} flag of the rest of the translation unit (the C++
    wrapper, the priority-queue traversal, etc.) is not on the critical path.
  \end{enumerate}

  The multi-target build still earns its keep on three counts: (i) it
  \textbf{guarantees correctness} on ARMv8.0 chips (no SIGILL on Kirin 659 /
  SD662 --- the original motivation),

  \begin{enumerate}
  \def\labelenumii{(\roman{enumii})}
  \setcounter{enumii}{1}
  \tightlist
  \item
    it leaves the architecture ready for future optimisations that \emph{do}
    require global flags (e.g.~an SVE path on ARMv9 chips, an i8mm path on
    Cortex-A78+), and
  \item
    the \texttt{+0.07\ MB} APK overhead is trivial. We document the null result
    here to pre-empt the reasonable reviewer question ``if the multi-target
    build is in the binary, why doesn't the dispatched p50 dominate the
    baseline-forced p50?''.
  \end{enumerate}

  \textbf{EMUI 10 packaging note}: collecting the FRL-L23 round 2
  baseline-forced data point required bypassing \texttt{adb\ install\ -r} (which
  refused with \texttt{DELETE\_FAILED\_INTERNAL\_ERROR} on five consecutive
  attempts across two reboot cycles and an \texttt{adb\ kill-server} reset) and
  using the multi-step
  \texttt{cmd\ package\ install-create\ /\ install-write\ /\ install-commit}
  flow directly through \texttt{adb\ shell}. The behaviour is documented Huawei
  iAware interaction with the package-installer service; the reproducer is in
  \texttt{experiment/storage/android/README.md}. No algorithmic finding is
  contingent on this workaround --- we report it for benchmarkers who hit the
  same lock.
\item
  \textbf{EMUI 10 on the Helio G80 unit (Y9a / FRL-L23) force-killed the bench
  process during the ObjectBox N = 20 000 ingest (≈ 12 min in native code).} A
  persistent foreground service plus a 5 s harness heartbeat plus an adb-side
  tickle was required to keep the process alive long enough to write the JSON.
  The cell is otherwise produced by the same harness as the reference run.
\end{enumerate}

The qualitative ordering (Dazzle SQ8 \textless{} ObjectBox 4.x ≪ SQLiteAI
precompute ≪ \texttt{sqlite-vec} ≪ raw SQLite) holds across all four chips and
across both ISA generations. The 95 \% CIs in the table above and in
\texttt{research/paper/vecbench\_cross\_platform\_ci.md} exclude 1.0 in every
chip × engine cell that the paper relies on.

\subsection{Algorithmic Asymmetry in the Vector
Benchmark}\label{algorithmic-asymmetry-in-the-vector-benchmark}

The comparison in §5.8 spans two algorithm classes: HNSW (Dazzle, ObjectBox) is
approximate nearest neighbour search with expected \(O(\log N)\) search
complexity \citep{hnsw}; SQLiteAI v0.9.95 implements SIMD-accelerated linear
scan \(O(N)\) over int8-quantised vectors. The asymmetry is by design of the
SQLiteAI product --- \texttt{vector\_quantize\_scan} is the only optimised query
path the product exposes; HNSW is not user-accessible in this version. Reporting
the speedup ratio against SQLiteAI without an algorithmic note is misleading, so
the §5.8.1 setup pins the disclosure inline.

Within Table 11 the operating-point reading is the one to use: \textbf{Dazzle
SQ8 lands at recall = 0.959 / p50 = 208 µs, ObjectBox 4.x at recall = 0.994 /
p50 = 853 µs, SQLiteAI precompute at recall = 0.987 / p50 = 1 407 µs}, all at N
= 20 000, dim = 384, recall floor \(\geq\) 0.95.

The two HNSW engines (Dazzle SQ8 and ObjectBox) are \textbf{two points on the
same recall--latency Pareto curve}, not winner / loser: Dazzle SQ8 trades int8
quantisation precision for \textasciitilde4× lower p50; ObjectBox keeps fp32
precision and pays the latency cost. \textbf{Both choices stay inside the
LLM-noise-floor envelope at this N} (208 µs and 853 µs are 0.007 \% and 0.030 \%
of a 2.84 s inference turn respectively; both are dwarfed by the inference
call). The recommendation is therefore not ``Dazzle SQ8 beats ObjectBox'': an
application that needs higher recall picks ObjectBox; one that needs the smaller
RAM footprint and is happy with int8 recall picks Dazzle SQ8; both are fine for
the agent-loop case. SQLiteAI sits in a different algorithm class (\(O(N)\)
scan); its 1 407 µs cell still fits the noise floor at N = 20 000 but not at N
\(\gg\) 20 000 --- the corpus-size threshold §5.8.4 documents.

Earlier drafts of this paper also reported a ``three operating points (recall
\(\geq\) 0.95, recall = 0.993, recall \(\geq\) 0.998)'' breakdown; the current
§5.8 runs report a single recall per engine because each engine sits at its
natural operating point given fixed bench parameters (M = 32, efC = 400). A
multi-recall sweep at fixed N is left to a follow-up Pareto-frontier study; the
direction would not change the envelope claim --- at any operating point in the
0.95--0.999 band, all three HNSW engines stay under the inference-noise-floor
threshold for typical mobile RAG corpora.

\section{Related Work}\label{related-work}

\textbf{On-device inference.} LiteRT-LM, llama.cpp, MLC-LLM, and ExecuTorch
\citep{litertlm, llamacpp, mlcllm, executorch} address model compression and
runtime efficiency. None addresses the stateful execution layer --- they treat
each inference call as an independent event.

\textbf{Memory for agents.} MemGPT \citep{memgpt} proposes a two-level memory
hierarchy (in-context + external) targeting server deployment. Generative Agents
\citep{generativeagents} uses a retrieval-augmented memory stream for NPC
simulation. A-MEM \citep{amem} introduces agent-managed memory with dynamic
linking. All assume cloud or server environments with a separate database
process, network, and substantial RAM. Dazzle moves the external memory layer
\textbf{inside the mobile app process}.

\textbf{Redis on constrained hardware.} Redis on Raspberry Pi achieves
sub-millisecond latency, but runs as a separate daemon over TCP. Valkey on
Android specifically has not been published to our knowledge.

\textbf{Embedded engines.} SQLite, LMDB, RocksDB are designed from the ground up
as in-process libraries. Dazzle takes the opposite path: starting from a server
engine designed around TCP loopback + single-threaded event loop and rewriting
its I/O substrate --- not its data structures --- to make in-process execution
the default. It sits between both camps: used as an embedded library, but
preserving the data model of a full server-grade database.

\textbf{Mobile vector databases --- three distinct products.} The mobile
SQLite-vector ecosystem includes three commonly-confused products that we
distinguish carefully here. (1) Plain \textbf{SQLite} has no native vector
support and is benchmarked only on its storage layer in §5.2. (2) Alex Garcia's
open-source \textbf{sqlite-vec} (\texttt{github.com/asg017/sqlite-vec}, Apache
2.0) is a SQLite extension adding brute-force vector search; this revision
benchmarks it in the SQLite-family sweep (§5.8.4). (3) \textbf{SQLiteAI
sqlite-vector} \citep{sqliteaivector}
(\texttt{github.com/sqliteai/sqlite-vector}, Elastic License 2.0) is the
commercial product from SQLiteAI Inc., which ships a SIMD-accelerated quantized
linear-scan query path (\texttt{vector\_quantize\_scan}) and deliberately avoids
HNSW. SQLiteAI sqlite-vector and the open-source sqlite-vec are both
commercially or publicly available as of 2026, and we benchmark SQLiteAI
directly in §5.8 because its production positioning makes it the most relevant
SQLite-ecosystem reference for an embedded-LLM-agent comparison; we make no
claim about it being the ``standard'' of the SQLite vector ecosystem.
\textbf{ObjectBox 4.x} \citep{objectbox4} is the other commercial peer in our
benchmark and ships native HNSW backed by float32 vectors.

\textbf{Server-grade vector ANN libraries.} \textbf{FAISS} \citep{faiss}
(Facebook AI), \textbf{DiskANN} \citep{diskann} (Microsoft Research), and
\textbf{Milvus} \citep{milvus} (Zilliz) set the open-source baseline for
approximate nearest-neighbour search at billion-point scale on server hardware.
The vector-quantisation literature underneath those engines --- Product
Quantisation \citep{jegoupq} in particular --- is the algorithmic precedent for
the SQ8 path Dazzle uses. None of FAISS, DiskANN, or Milvus is a mobile
deployment target --- FAISS expects x86 + libomp, DiskANN's SSD-resident graph
assumes a workstation-class IO subsystem, Milvus is a distributed system with a
separate orchestrator --- but all three are the algorithmic precedent against
which any HNSW implementation (including Dazzle's, ObjectBox's, and
\texttt{vector\_quantize\_scan}) is compared at the index-construction level.
Dazzle's NEON kernels for fp32 dot product, fp32 L2-squared, fp16 dot product,
and i8 cosine come from \texttt{simsimd} \citep{simsimd} (Vardanian, MIT
licence) and use that library's runtime dispatcher
(\texttt{simsimd\_capabilities()} over \texttt{/proc/cpuinfo}) to select the
best available variant on each chip, which is the mechanism §6.3 exercises. The
contribution of this paper is not a new ANN algorithm; it is putting an HNSW
index a \texttt{getResources()} reference away from a Kotlin/Swift app on a
budget Android phone, on the same primitive surface as five non-vector
data-structure commands.

\textbf{Mobile object stores beyond SQLite.} \textbf{Realm} \citep{realm} (now
MongoDB) and \textbf{WCDB} \citep{wcdb} (Tencent / WeChat) are the two
production-deployed peers to ObjectBox in the mobile object-store space. Both
target Android + iOS, both ship higher-level typed APIs than raw SQLite, and
both are widely used in commercial apps. Neither offers vector search or
HyperLogLog as a primitive (Realm ships full-text search; WCDB is
SQLite-on-the-wire underneath), so they sit in the same category as ObjectBox
--- embedded object stores with two to four primitive types --- and would extend
the ``≤ 4 primitives'' column of Table 1 rather than challenge it.

\textbf{RAG evaluation.} \textbf{RAGTruth} \citep{ragtruth} is the canonical
hallucination corpus for retrieval-augmented LLM evaluation as of 2024. Section
§5.9 deliberately uses NQ-open instead because the unit of measurement here is
\emph{recall of canonical short-answer strings} (whether the LLM emits the gold
span when retrieval is present); RAGTruth measures hallucination \emph{given}
retrieved context, which is a strictly downstream signal that requires a
hallucination judge. Layering RAGTruth-style evaluation on top of the §5.9 setup
is future work and would belong on the same axis as the EM-vs-F1 reporting we
already include.

\section{Conclusion}\label{conclusion}

Dazzle argues for a different backend-selection lens for on-device agents:
primitive surface, developer ergonomics, and latency envelope matter more than
isolated microsecond races. In the measured workload, all evaluated backends
operate inside an acceptable retrieval envelope; the operational differentiators
are which agent-state primitives are available natively and what engineering
effort is required to use them.

Three empirical results support that framing. First, retrieval latency in
Dazzle's steady-state path is roughly 0.00176 \% of a single on-device LLM
inference turn (50 µs / \textasciitilde2.84 s for Gemma 4 E2B on the Moto G35
5G), placing backend retrieval cost within the noise floor of the end-to-end
loop. Second, primitive surface and development effort differ materially: Dazzle
exposes ten native primitives versus \(\leq\) 4 in the measured embedded
alternatives, with representative agent-state implementations around 175 LOC
(Android) / 186 LOC (iOS) in Dazzle versus roughly 200--290 LOC across the
alternatives (Table 9, both platforms). Third, at the application level, a 2×2
RAG ablation lifts \texttt{EM\_contains} 6.0× on Qwen 2.5 0.5B (0.105 → 0.630)
and 6.7× on Qwen 2.5 1.5B (0.110 → 0.735); a 380\textasciitilde MB model with
retrieval beats a 3× larger model without retrieval on every factual metric,
end-to-end on a \$150-class Android device.

The practical question for the next on-device deployment therefore shifts from
``which backend is fastest?'' to ``which primitive set and which model
composition best fit the application, given that retrieval cost is effectively
invisible at turn level.''

Dazzle is publicly distributed as \texttt{dazzle-sdk} 1.0.0-beta.4 across Maven
Central, pub.dev, npm, and Swift Package Manager. Dazzle-origin code is
Apache-2.0; derivative Valkey portions retain BSD-3-Clause. We invite the
community to validate these results on additional mobile hardware and extend the
benchmark matrix across broader workloads.

\subsection{Artifact and Reproducibility}\label{artifact-and-reproducibility}

To make the evaluation auditable, this repository includes the scripts, dataset
builders, and raw result trees used to produce the paper tables. The artifact
package is organized as follows:

\begin{itemize}
\tightlist
\item
  Benchmark runners: \texttt{research/scripts/run\_experiment.sh},
  \texttt{research/scripts/run\_full\_benchmark.sh},
  \texttt{research/scripts/run\_ios\_benchmark.sh},
  \texttt{research/scripts/run\_ablation\_sweep.sh},
  \texttt{research/scripts/run\_storage\_microbench\_per\_backend.sh},
  \texttt{research/scripts/run\_vector\_sqlite\_family.sh}.
\item
  Dataset generation and RAG slice tooling:
  \texttt{research/scripts/generate\_dataset.py},
  \texttt{research/scripts/nq\_slice.py},
  \texttt{research/scripts/recompute\_rag\_metrics.py}.
\item
  Result analysis and table generation:
  \texttt{research/scripts/analyze\_results.py},
  \texttt{research/scripts/analyze\_storage\_microbench.py},
  \texttt{research/scripts/make\_vector\_bench\_table.py},
  \texttt{research/scripts/analyze\_vector\_sqlite\_family.py}.
\item
  Raw and processed result trees: \texttt{research/benchmarks/results/},
  \texttt{research/results/}.
\item
  Input data snapshots: \texttt{research/data/}, including
  \texttt{research/data/nq\_slice/}.
\end{itemize}

All headline metrics in §5 were computed from these local artifacts. A compact
step-by-step reproduction guide is provided in
\texttt{research/paper/REPRODUCIBILITY.md} and can be executed end-to-end on the
same workspace layout.

\subsection{Versioning Policy and Public Commitment to
Update}\label{versioning-policy-and-public-commitment-to-update}

This paper documents a moving target: \texttt{dazzle-sdk} is a live artifact
under active development across four package registries (Maven Central, pub.dev,
npm, Swift Package Manager), the comparison backends evolve (e.g., a future
SQLiteAI release with user-accessible HNSW would shift §5.8 from ``HNSW vs
brute-force scan'' to ``HNSW vs HNSW''), and several findings in this revision
are explicitly null results that invite further investigation (the Cortex-A75
codegen regression in §6.3 footnote ², the apples-to- apples vs dispatched
parity on Unisoc T760, the marginal \texttt{simsimd} runtime-dispatch overhead).
To keep the paper a useful reference rather than a frozen snapshot, the author
commits to the following versioning policy on arXiv:

\begin{enumerate}
\def\labelenumi{\arabic{enumi}.}
\tightlist
\item
  \textbf{Each measurable change in any headline cell of Tables 3, 5, 11, 12, or
  15 produces a new arXiv revision.} ``Measurable'' is operationalised as: the
  new 95 \% bootstrap CI for the cell does not overlap the previous revision's
  CI on the same operating point, on the same physical device. Pure noise
  re-runs do not trigger a revision.
\item
  \textbf{Each new physical device added to the cross-platform table (§6.3)
  produces a new arXiv revision.} The current revision covers four Android SoCs
  (Unisoc T760, MediaTek Helio G80, Snapdragon 662, Kirin 659) plus iPhone 12
  Pro on the storage-only path; chips and OS versions added later --- Snapdragon
  8-series, Apple A17, Tensor G3, Cortex-X1+ class, future ARMv9 cores --- land
  as a v2 / v3 / \ldots{} with the new rows appended in-place.
\item
  \textbf{Each accepted external pull request that improves a non-Dazzle
  backend's measured numbers in Tables 3 / 5 / 11 produces a new arXiv revision}
  --- with the contributing PR cited and the old numbers preserved as a
  strikethrough in the same cell so a reader can see what changed and when. This
  is the operational half of the conflict-of-interest commitment listed in §6.3.
\item
  \textbf{The arXiv abstract carries a ``Last updated'' date and a one-line
  changelog of the most recent material change} so a citation to a specific
  revision points to a specific set of numbers.
\item
  \textbf{All raw JSON measurements} for every cell in every revision live in
  \href{https://github.com/IvanAliaga/dazzle/tree/main/research/benchmarks/results}{\texttt{research/benchmarks/results/}}
  under timestamped paths, never overwritten across revisions. Bootstrap-CI
  re-derivations (\texttt{research/scripts/\ bootstrap\_*\_lats.py}) are
  deterministic (\texttt{B\ =\ 10\ 000}, \texttt{seed\ =\ 42}); a reader at any
  point in time can re-derive every CI in the table from the timestamped JSONs
  alone.
\item
  \textbf{The repository main branch is the source of truth.} When a discrepancy
  arises between an arXiv revision and \texttt{main}, \texttt{main} wins; the
  discrepancy triggers a revision per (1) above. The commit hash of the head of
  \texttt{main} at the moment of each arXiv submission is recorded in §8.1 of
  that revision.
\end{enumerate}

Out of scope for the versioning policy: typo fixes, prose tweaks, citation
reformats, and minor restructuring without changes to measurements --- those are
batched and produce a revision only when they accumulate alongside a
measurement-driven trigger.

The author is reachable at \texttt{ivan.aliaga@urp.edu.pe} for reviewer
questions, reproduction issues, and pull requests against the benchmark code.
Issues opened on the GitHub repository are acknowledged within seven days;
numbered arXiv revisions are processed within fourteen days of the trigger
event. This is a public commitment, not best-effort phrasing --- failures to
meet it are corrigible by the same channels.

\section{References}\label{references}

\section{Appendix A --- LLM Adapter Layer
(Reference)}\label{appendix-a-llm-adapter-layer-reference}

The Dazzle SDK exposes a single \texttt{LLMClient} interface (completion,
streaming, tool-call deltas, lifecycle) and ships five concrete adapters behind
it: \texttt{LlamaCppClient} and \texttt{LiteRtLmClient} (on-device GGUF /
LiteRT-LM runtimes), \texttt{FoundationModelsClient} (Apple Foundation Models),
\texttt{OpenAICompatibleClient} (OpenAI-format chat + tool-calling, cloud or
local), and \texttt{AnthropicClient} (Anthropic Messages API + streaming content
blocks). The interface exists for the same reason the storage layer exists: the
agent loop should not change when the underlying provider does.

The full engineering reference for the adapter layer --- wire-format
normalisation across providers, the unified \texttt{Delta} event type
(\texttt{text} / \texttt{toolUseStart} / \texttt{toolUseFinish} / \texttt{usage}
/ \texttt{end}), the three Flutter \texttt{EventChannel} invariants applied
preventively to every bridge, and the live-verification matrix against the
production Anthropic Messages API on four cross-stack configurations --- lives
in the SDK technical report at \texttt{docs/sdk/llm\_adapters.md} (cited from §3
of the companion engineering report). It is documentation of the SDK surface,
not a contribution of this paper, and is kept out of the main body to avoid
bloating the manuscript with material that does not bear on the storage thesis.


% --- bibliography ---
\bibliographystyle{unsrt}
\bibliography{refs}

\end{document}
